\documentclass[12pt,a4paper]{article}

\usepackage{fontspec}
\setmainfont{Times New Roman}
\usepackage{amsmath,amssymb,bm}
\usepackage[top=2.5cm,bottom=2.5cm,left=3cm,right=2.5cm]{geometry}
\usepackage{setspace}
\onehalfspacing
\usepackage{indentfirst}
\setlength{\parindent}{1.5em}
\setlength{\parskip}{0.35em}
\setlength{\emergencystretch}{1em}
\usepackage{graphicx}
\graphicspath{{figures/}}
\usepackage{booktabs}
\usepackage{multirow}
\usepackage{array}
\usepackage{tabularx}
\usepackage{longtable}
\usepackage{xcolor}
\usepackage{float}
\usepackage{caption}
\usepackage{subcaption}
\usepackage{tikz}
\usetikzlibrary{arrows.meta,positioning,fit,calc,shapes.geometric}
\usepackage{adjustbox}
\usepackage{pdflscape}
\usepackage{hyperref}
\usepackage{url}
\usepackage{cite}
\newcolumntype{P}[1]{>{\raggedright\arraybackslash}p{#1}}
\newcolumntype{Y}{>{\raggedright\arraybackslash}X}

\hypersetup{
  colorlinks=true,
  linkcolor=blue!55!black,
  citecolor=blue!55!black,
  urlcolor=blue!55!black,
  pdftitle={Phân loại số người đóng góp trong hỗn hợp DNA và Nhận dạng cá nhân trong hỗn hợp DNA}
  pdfauthor={Nguyễn Vũ Mạnh, Lê Sỹ Nguyên, Nguyễn Đắc Quang, Trịnh Hoàng Giang, Nguyễn Quốc Dũng}
}

\renewcommand\thesection{\Roman{section}}
\renewcommand\thesubsection{\arabic{section}.\arabic{subsection}}
\renewcommand\thesubsubsection{\arabic{section}.\arabic{subsection}.\arabic{subsubsection}}
\renewcommand{\contentsname}{Mục lục}
\renewcommand{\figurename}{Hình}
\renewcommand{\tablename}{Bảng}
\renewcommand{\refname}{Tài liệu tham khảo}

\definecolor{peakblue}{RGB}{43,108,176}
\definecolor{locusgreen}{RGB}{47,133,90}
\definecolor{profileorange}{RGB}{221,107,32}
\definecolor{headpurple}{RGB}{128,90,213}
\definecolor{lightgray}{RGB}{245,247,250}

\tikzset{
  block/.style={draw,rounded corners=2pt,align=center,minimum height=1.15cm,
    text width=2.35cm,fill=blue!6,font=\small},
  smallblock/.style={draw,rounded corners=2pt,align=center,minimum height=0.95cm,
    text width=1.85cm,fill=gray!7,font=\scriptsize},
  line/.style={-{Latex[length=2.4mm]},thick},
  groupbox/.style={draw,dashed,rounded corners=4pt,inner sep=7pt}
}

\begin{document}

\begin{titlepage}
  \centering
  \vspace*{1.0cm}

  {\large\textbf{TRƯỜNG ĐẠI HỌC BÁCH KHOA HÀ NỘI}}\\[0.25cm]
  {\normalsize Trường Công nghệ Thông tin và Truyền thông}\\[1.6cm]

  \rule{\linewidth}{1.4pt}\\[0.5cm]
  {\LARGE\bfseries Phân loại số người đóng góp trong hỗn hợp DNA}\\[0.25cm]
  {\LARGE\bfseries và Nhận dạng cá nhân trong hỗn hợp DNA}\\[0.5cm]
  \rule{\linewidth}{1.4pt}\\[1.4cm]

  {\large\textbf{Báo cáo nghiên cứu}}\\[1.3cm]

  \begin{tabular}{>{\bfseries}rl}
    Thành viên nhóm: & Nguyễn Vũ Mạnh \\[0.25cm]
                     & Lê Sỹ Nguyên \\[0.25cm]
                     & Nguyễn Đắc Quang \\[0.25cm]
                     & Trịnh Hoàng Giang \\[0.25cm]
                     & Nguyễn Quốc Dũng
  \end{tabular}

  \vfill
  {\large Hà Nội, tháng 6 năm 2026}
\end{titlepage}

\tableofcontents
\newpage

\section*{Tóm tắt}
\addcontentsline{toc}{section}{Tóm tắt}

Phân tích hỗn hợp DNA là một trong những nhiệm vụ khó nhất của giám định di truyền pháp y vì tín hiệu điện di quan sát được có thể đồng thời chứa allele của nhiều người, stutter, nhiễu nền, allele dropout và chênh lệch lớn về tỷ lệ đóng góp. Hai câu hỏi xuất hiện sớm trong quy trình diễn giải là mẫu có bao nhiêu người đóng góp và những người nào trong một tập tham chiếu đã góp DNA vào mẫu. Nghiên cứu này trình bày DS-ST như một họ mô hình dùng chung cách biểu diễn hồ sơ STR dưới dạng tập peak bất biến đối với hoán vị. Đối với Number of Contributors, mỗi peak được chuyển thành token tám trường gồm chỉ số locus, allele, log chiều cao, tỷ lệ chiều cao trong locus, tỷ số với peak cha stutter, thứ hạng chiều cao, số peak trong locus và chiều cao tương đối toàn hồ sơ. Bảy trường số liên tục được chuẩn hóa và mã hóa bằng periodic numerical embedding, còn locus được mã hóa bằng embedding học được. Chuỗi token đi qua hai khối ISAB++ có 32 inducing point, bốn attention head, SetNorm, clean path residual và non competitive sigmoid attention tại bước inducing point thu nhận peak. Pooling by Multihead Attention tạo véc tơ hồ sơ 128 chiều; đầu NOC $128\!\to\!64\!\to\!5$ được fine-tune từ checkpoint tiền huấn luyện trên nhiệm vụ contributor identification. Đối với nhận dạng contributor, cùng họ backbone được mở rộng bằng feasibility filter phân tách peak riêng và peak chia sẻ, cùng bộ giải mã Adaptive Slot gồm khởi tạo slot từ genotype theo Conditional Slot Attention \cite{ma2024cosa}, tinh chỉnh attribution theo Guided Slot Attention \cite{lee2024gsa}, phân phối peak đến donor qua Sinkhorn optimal transport theo MESH \cite{zhang2023mesh} và gate tồn tại theo Adaptive Slot Attention \cite{fan2024adaslot}. Nhánh này còn có đầu từ chối open set, đầu attribution theo peak, hồi quy mixture proportion và bước phi reranking sau mô hình. Dữ liệu NoC được xây dựng từ PROVEDIt, cơ sở dữ liệu công khai có hơn 27.000 hỗn hợp DNA pháp y \cite{provedit2018,proveditweb}. Tập GlobalFiler gồm 10.124 hồ sơ thuộc 3.418 nhóm hỗn hợp vật lý; nhãn một người chiếm 8.119 hồ sơ, vì vậy quá trình huấn luyện sử dụng oversampling, trọng số lớp và Focal Loss. Trên GroupKFold năm phần với 10.124 hồ sơ GlobalFiler, DS-ST đạt Macro F1 $0{,}9444\pm0{,}0042$, vượt TabPFN ($0{,}9297\pm0{,}0108$), XGBoost ($0{,}9230\pm0{,}0209$), CNN một chiều ($0{,}8854\pm0{,}0100$) và Forensic Transformer ($0{,}8846\pm0{,}0131$). Ablation ba phần cho thấy bỏ tiền huấn luyện làm Macro F1 giảm từ $0{,}9122$ xuống $0{,}8918$, còn pipeline thế hệ trước với token ba trường, embedding thô và ISAB cơ sở chỉ đạt $0{,}8430$. Các kết quả này cho thấy tiền huấn luyện từ nhiệm vụ contributor, biểu diễn peak giàu thông tin và attention dành riêng cho tập hợp đều có đóng góp thực nghiệm đáng kể.

\medskip
\noindent\textbf{Từ khóa:} hỗn hợp DNA; số lượng người đóng góp; STR; DS-ST; Set Transformer; học sâu pháp y; giải trộn DNA; PROVEDIt.

\section*{Các đóng góp chính}
\addcontentsline{toc}{section}{Các đóng góp chính}

Đóng góp thứ nhất của nghiên cứu là xây dựng một biểu diễn thống nhất cho hai bài toán Number of Contributors và contributor identification, trong đó cùng một backbone biến đổi tín hiệu peak thành các biểu diễn peak, locus và toàn hồ sơ. Thiết kế này tránh việc xây dựng hai pipeline tiền xử lý tách rời, đồng thời bảo toàn tính bất biến đối với hoán vị peak, vốn là yêu cầu tự nhiên của dữ liệu dạng tập hợp \cite{zaheer2017,lee2019}.

Đóng góp thứ hai là kiến trúc DS-ST nhánh NoC cho bài toán dự đoán Number of Contributors. Mô hình sử dụng token tám trường mã hóa từng peak, Periodic Linear Representation \cite{gorishniy2022} cho các đặc trưng số liên tục, hai khối ISAB++ với sigmoid attention không cạnh tranh, SetNorm, clean-path residual, 32 inducing point và Pooling by Multihead Attention. Backbone được tiền huấn luyện trên nhiệm vụ contributor identification 45 lớp rồi fine-tune cho NOC bằng quy trình hai giai đoạn: đóng băng backbone năm epoch warmup, sau đó mở toàn mạng với Focal Loss và trọng số lớp nghịch đảo tần suất.

Đóng góp thứ ba là xây dựng nhánh xác định contributor với 45 donor tham chiếu bằng bộ giải mã slot có điều kiện theo genotype. Các khái niệm Conditional Slot Attention, Guided Slot Attention, MESH và Adaptive Slot Attention được điều chỉnh sang bối cảnh pháp y, trong đó mỗi slot đại diện cho một donor, ma trận Sinkhorn biểu diễn phân bổ bằng chứng peak đến donor và gate Gumbel Sigmoid biểu diễn khả năng donor tồn tại trong hỗn hợp \cite{ma2024cosa,lee2024gsa,zhang2023mesh,fan2024adaslot}.

Đóng góp thứ tư là thiết kế quy trình đánh giá chống rò rỉ theo danh tính hỗn hợp vật lý, sử dụng GroupKFold để các lần tiêm hoặc thời gian thu nhận khác nhau của cùng một mixture không xuất hiện đồng thời trong tập huấn luyện và kiểm định. Nghiên cứu đồng thời báo cáo benchmark với XGBoost, CNN, Forensic Transformer, TabPFN và DS-ST; phân tích ablation về tiền huấn luyện và enrichment; đồng thời xây dựng dữ liệu cho GlobalFiler, Identifiler Plus và PowerPlex 16 để chuẩn bị đánh giá khác miền.

\newpage
\section*{Danh mục từ viết tắt}
\addcontentsline{toc}{section}{Danh mục từ viết tắt}

\begin{longtable}{@{}P{3.5cm}P{11.5cm}@{}}
AdaSlot   & Adaptive Slot Attention \\
CE        & Capillary Electrophoresis (điện di mao quản) \\
CNN       & Convolutional Neural Network (mạng nơ-ron tích chập) \\
CoSA      & Conditional Slot Attention \\
DNA       & Deoxyribonucleic Acid \\
DS-ST     & Donor-Slot Set Transformer \\
GF        & GlobalFiler (bộ kit giám định pháp y) \\
GSANet    & Guided Slot Attention Network \\
ISAB      & Induced Set Attention Block (khối attention tập hợp quy nạp) \\
ISAB++    & Enhanced Induced Set Attention Block (ISAB cải tiến) \\
MAB       & Multihead Attention Block \\
MESH      & Matching via Entropy-regularized Sinkhorn (optimal transport) \\
ML        & Machine Learning (máy học) \\
MLP       & Multilayer Perceptron (perceptron đa lớp) \\
MPS       & Massively Parallel Sequencing (giải trình tự song song) \\
NOC       & Number of Contributors (số người đóng góp DNA) \\
OT        & Optimal Transport (vận chuyển tối ưu) \\
PACE      & Probabilistic Assessment for Contributor Estimation \\
PCR       & Polymerase Chain Reaction (phản ứng chuỗi polymerase) \\
PLR       & Periodic Linear Representation \\
PMA       & Pooling by Multihead Attention \\
PROVEDIt  & Project Research Openness for Validation with Empirical Data \\
RF        & Random Forest (rừng ngẫu nhiên) \\
STR       & Short Tandem Repeat (lặp đơn vị ngắn tandem) \\
TAWSEEM   & Tool for Automated Weighted Scoring and Estimation of Evidence for Mixtures \\
XAI       & Explainable Artificial Intelligence (AI có khả năng giải thích) \\
XGBoost   & Extreme Gradient Boosting \\
\end{longtable}

\newpage
\section{Giới thiệu}

\subsection{Bối cảnh giám định DNA hỗn hợp}

DNA profiling dựa trên short tandem repeat đã trở thành một trụ cột của giám định sinh học kể từ khi các mẫu lặp đa hình được sử dụng cho nhận dạng cá thể vào thập niên 1980 \cite{jeffreys1985}. Trong quy trình capillary electrophoresis hiện đại, các vùng STR được khuếch đại, phân tách theo kích thước và biểu diễn thành các peak huỳnh quang. Một hồ sơ nguồn đơn thường có không quá hai allele thật tại mỗi locus autosomal, nhưng hồ sơ hỗn hợp có thể chứa nhiều hơn hai allele, các peak có độ cao chênh lệch mạnh và các artefact phát sinh từ phản ứng PCR hoặc thiết bị. PROVEDIt được xây dựng để hỗ trợ kiểm chứng các phương pháp phân tích trong bối cảnh này và hiện được Rutgers mô tả là cơ sở dữ liệu mở gồm hơn 27.000 hỗn hợp DNA có ý nghĩa pháp y \cite{provedit2018,proveditweb}.

Trong một vụ việc thực tế, mẫu DNA có thể được thu từ máu, tinh dịch, tế bào biểu mô hoặc dấu vết tiếp xúc. Số người để lại vật chất sinh học thường không được biết trước. Việc ước lượng Number of Contributors vì thế là bước đầu của nhiều quy trình diễn giải. NOCIt, STRmix, EuroForMix và các hệ thống probabilistic genotyping khác đều cần mô hình hóa một hoặc nhiều giả thuyết về số contributor; giả thuyết sai có thể làm thay đổi likelihood ratio và cách đánh giá sức nặng của chứng cứ \cite{haned2011,bleka2016,bright2016,coble2019}. NIST nhấn mạnh rằng khả năng phát hiện DNA từ lượng tế bào ngày càng nhỏ làm tăng số trường hợp xuất hiện hỗn hợp phức tạp và khiến việc phân biệt nguồn trở nên khó hơn \cite{nist2024}.

Phương pháp đơn giản nhất thường dựa trên maximum allele count. Nếu locus có \(A_{\max}\) allele, một cận dưới tự nhiên cho số người là \(\lceil A_{\max}/2\rceil\). Quy tắc này dễ giải thích nhưng mất thông tin về chiều cao peak, tần suất allele, dropout và stutter. Haned và cộng sự đã xây dựng cách tiếp cận xác suất so sánh giả thuyết hai người, ba người và các số contributor khác bằng xác suất allele quan sát được \cite{haned2011}. Alfonse và cộng sự đánh giá NOCIt, maximum likelihood và maximum allele count trên 100 mixture phức tạp, cho thấy việc suy luận trở nên khó hơn rõ rệt ở chế độ low template \cite{alfonse2017}. Những kết quả này cho thấy NOC không chỉ là bài toán đếm số allele mà là bài toán suy luận từ tín hiệu nhiễu, thiếu và chồng lấn.

\subsection{Hai bài toán liên kết nhưng không đồng nhất}

Bài toán thứ nhất của nghiên cứu là dự đoán \(y^{\mathrm{noc}}\in\{1,2,3,4,5\}\) từ một hồ sơ STR. Đây là bài toán phân loại năm lớp có thứ tự, nhưng notebook cuối cùng giải quyết bằng đầu softmax năm lớp vì mục tiêu thực nghiệm là tối ưu Macro F1 dưới phân phối mất cân bằng. Mã nguồn contributor từng khảo sát CORN và các biến thể đếm từ slot mass, nhưng README của pipeline sạch ghi nhận các đầu này làm suy giảm mạnh lớp NOC ba và bốn; vì vậy chúng bị loại khỏi kiến trúc contributor cuối cùng. Đối với nhánh NoC độc lập, đầu được sử dụng thực tế là MLP \(128\rightarrow64\rightarrow5\), huấn luyện bằng Focal Loss và class weight. Việc phân biệt rõ thành phần đã dùng với thành phần từng được khảo sát giúp tránh mô tả một kiến trúc không trùng với notebook.

Bài toán thứ hai là xác định tập contributor từ một panel tham chiếu có \(C\) donor. Nhãn của một hồ sơ là véc tơ đa nhãn \(\bm{y}^{\mathrm{id}}\in\{0,1\}^{C}\), trong đó số phần tử bằng một chính là NOC thật nếu mọi contributor đều nằm trong panel. Bài toán này khó hơn phân loại NOC vì mô hình phải phân biệt các donor chia sẻ allele và phải tìm bằng chứng của contributor thiểu số khi peak của họ bị che bởi contributor chính. Nghiên cứu sequencing của Phan và cộng sự sử dụng 27 STR và 94 SNP từ sáu cá thể để dự đoán cá thể trong hỗn hợp, minh họa tiềm năng của học sâu đối với nhận dạng contributor nhưng cũng cho thấy thiết lập panel nhỏ và dữ liệu MPS khác đáng kể so với capillary electrophoresis \cite{phan2021}.

Hai bài toán có thể hỗ trợ nhau nhưng không nên bị đồng nhất. NOC cung cấp cardinality cho phép giải mã top \(k\) ở nhánh nhận dạng, trong khi phân bố xác suất donor và tỷ lệ hỗn hợp có thể cung cấp tín hiệu ngược cho đầu đếm. Tuy nhiên, một mô hình NOC có accuracy cao chưa chắc xác định đúng danh tính, và một mô hình xếp hạng donor tốt vẫn có thể chọn sai số lượng nếu cardinality không được hiệu chỉnh. Vì vậy, khung đề xuất sử dụng backbone chung nhưng giữ loss và metric riêng cho từng nhiệm vụ.

\subsection{Động lực của biểu diễn tập hợp phân cấp}

Véc tơ phẳng theo cột allele có ưu điểm đơn giản và phù hợp với XGBoost hoặc TabPFN, nhưng cấu trúc của nó phụ thuộc bộ kit và không biểu diễn trực tiếp quan hệ giữa các peak trong cùng locus. CNN một chiều áp dụng convolution dọc theo véc tơ cột, song tính lân cận giữa hai cột không nhất thiết tương ứng với quan hệ sinh học. Forensic Transformer cải thiện điều này bằng cách tạo một token cho mỗi locus, nhưng việc nén toàn bộ peak của một locus thành một véc tơ trước attention có thể làm mất bằng chứng yếu. Dữ liệu STR phù hợp tự nhiên hơn với cấu trúc tập hợp lồng nhau, trong đó hồ sơ là tập các locus và mỗi locus là tập peak.

Deep Sets chứng minh rằng một lớp hàm bất biến hoán vị có thể được biểu diễn dưới dạng \(\rho(\sum_i\phi(\bm{x}_i))\) trong các điều kiện thích hợp \cite{zaheer2017}. Set Transformer mở rộng ý tưởng này bằng attention, cho phép mô hình hóa tương tác giữa các phần tử trong tập và dùng inducing points để giảm độ phức tạp từ bậc hai xuống xấp xỉ \(O(NM)\), với \(N\) phần tử và \(M\) inducing points \cite{lee2019}. Kiến trúc thực tế của đồ án đi theo hướng thứ hai: mỗi peak là một phần tử của tập, chỉ số locus được đưa vào embedding, còn quan hệ giữa mọi peak được học qua ISAB++. Thiết kế không dựng tensor cố định theo từng locus và không dùng một Cross Locus Transformer riêng; cấu trúc locus được bảo toàn trong chính token và trong bảng tra genotype.

\section{Các nghiên cứu liên quan}

Lĩnh vực ước lượng số người đóng góp DNA hỗn hợp đã trải qua ba thế hệ phương pháp kế tiếp nhau: từ quy tắc đếm allele dựa trên luật, qua machine learning với đặc trưng thủ công, đến deep learning trên dữ liệu quy mô lớn hoặc mô phỏng. Bảng~\ref{tab:related_summary} tổng hợp ba công trình tiêu biểu đại diện cho ba thế hệ đó; các phần tiếp theo sẽ trình bày chi tiết nội dung, phương pháp và đóng góp của từng nhóm công trình.

\begin{table}[H]
\centering
\caption{Tổng hợp ba nghiên cứu tiêu biểu về ước lượng Number of Contributors trong hỗn hợp DNA STR.}
\label{tab:related_summary}
\adjustbox{max width=\linewidth}{%
\small
\begin{tabular}{@{}P{3.2cm}lP{3.6cm}P{5.2cm}P{5.2cm}@{}}
\toprule
\textbf{Nghiên cứu} & \textbf{Năm} & \textbf{Phương pháp chính} & \textbf{Điểm mạnh} & \textbf{Khoảng trống / Hạn chế} \\
\midrule
Marciano \& Adelman (PACE) \cite{marciano2017}
  & 2017
  & So sánh sáu thuật toán ML (SVM, MLP, LR, kNN, DT) trên đặc trưng thống kê STR; lựa chọn đặc trưng bằng KL divergence; phân loại dưới 1 giây
  & Không dùng dữ liệu mô phỏng; ${\sim}97\%$ accuracy trên tập kiểm thử thực; khung so sánh đa thuật toán rõ ràng; thực thi thời gian thực
  & Chỉ bốn lớp NOC (1–4 người); bộ kit hẹp (1.405 profile Identifiler, 20 cá thể); đặc trưng thủ công không mô hình hóa cấu trúc nội locus; khó chuyển sang bộ kit GlobalFiler hay PROVEDIt \\
\addlinespace
Alamoudi et al.\ (TAWSEEM) \cite{alamoudi2022}
  & 2022
  & MLP sâu trên véc tơ phẳng ${\sim}$6.000 chiều bao phủ bốn multiplex STR từ PROVEDIt; đánh giá ba kịch bản tập dữ liệu khác nhau về quy mô
  & $97\%$ accuracy trong kịch bản bốn multiplex với 1–5 người; huấn luyện trên PROVEDIt quy mô lớn; cho thấy MLP sâu vượt random forest trên dữ liệu lớn
  & Véc tơ đầu vào phụ thuộc chặt vào bộ kit và số allele bin; phương pháp phân chia dữ liệu không ngăn rò rỉ mixture identity (không GroupKFold); chưa công bố mã nguồn để tái lập kết quả \\
\addlinespace
Taylor \& Humphries (deepNoC) \cite{taylor2024deepnoc}
  & 2024
  & Mạng sâu DNN huấn luyện trên 100.000 profile STR được mô phỏng; fine-tune thêm bằng profile casework thực của phòng thí nghiệm; đầu ra phụ hỗ trợ giải thích
  & Mở rộng NOC đến 10 người đóng góp; chiến lược transfer learning từ mô phỏng sang thực nghiệm; cho phép thích nghi với điều kiện riêng từng phòng lab
  & Chất lượng mô phỏng quyết định chất lượng tiền huấn luyện; khoảng cách giữa synthetic signal và casework thật chưa được định lượng; mô hình không khai thác cấu trúc tập peak theo locus \\
\bottomrule
\end{tabular}%
}
\end{table}

\subsection{Phương pháp thống kê và quy tắc đếm allele}

Maximum allele count là cận dưới dễ tính cho NOC, nhưng hiệu năng suy giảm khi các contributor chia sẻ allele hoặc khi dropout làm allele thật biến mất. Haned và cộng sự đề xuất tính likelihood cho các giả thuyết số contributor bằng cả số lượng, độ lớn và tần suất allele, nhờ đó biến NOC thành bài toán so sánh mô hình xác suất \cite{haned2011}. Alfonse và cộng sự sau đó so sánh NOCIt, maximum likelihood estimator và maximum allele count trên 100 mixture phức tạp, tập trung vào giới hạn low template \cite{alfonse2017}. NOCIt được kiểm chứng quy mô lớn trên 815 profile GlobalFiler gồm 100 mẫu một người, 193 mẫu hai người, 170 mẫu ba người, 186 mẫu bốn người và 166 mẫu năm người trong một nghiên cứu validation \cite{grgicak2020}. Điểm mạnh của các phương pháp xác suất là có thể liên hệ trực tiếp với likelihood ratio, nhưng chi phí tính toán, giả định mô hình và nhu cầu cấu hình chuyên gia khiến chúng không phải lúc nào cũng phù hợp cho sàng lọc nhanh.

EuroForMix sử dụng mô hình liên tục để đánh giá hồ sơ STR có artefact và đã trở thành một tham chiếu quan trọng trong probabilistic genotyping nguồn mở \cite{bleka2016}. STRmix cũng được developmental validation cho diễn giải mixture và được sử dụng rộng rãi trong thực hành \cite{bright2016}. Các hệ thống này giải quyết bài toán đánh giá chứng cứ toàn diện hơn một classifier NOC, nhưng NOC thường vẫn phải được đặt trước hoặc đánh giá qua nhiều giả thuyết. Do đó, một mô hình dự đoán NOC nhanh không thay thế probabilistic genotyping mà đóng vai trò đề xuất giả thuyết và kiểm tra tính hợp lý trước khi chạy mô hình xác suất.

\subsection{Machine learning cổ điển cho Number of Contributors}

PACE của Marciano và Adelman là một trong những hệ thống machine learning sớm cho NOC. Công trình sử dụng 1.405 profile không mô phỏng tạo từ 20 cá thể, một đến bốn contributor, 15 locus Identifiler, và đánh giá sáu thuật toán gồm nearest neighbor, decision tree, logistic regression, multilayer logistic regression, multilayer perceptron và support vector machine. Kullback Leibler divergence được dùng để loại đặc trưng có divergence dưới 0,01. Hệ thống báo cáo khoảng 98\% accuracy ở training và 97\% ở testing, đồng thời thực hiện phân loại trong thời gian dưới một giây \cite{marciano2017,alamoudi2022}. Hạn chế chính là dữ liệu và bộ kit hẹp hơn PROVEDIt, trong khi feature engineering được thiết kế riêng cho pipeline.

Benschop và cộng sự xây dựng tập 590 PowerPlex Fusion 6C profile từ 1.174 donor, bao phủ một đến năm contributor và các mức allele sharing, dropout, mixture proportion cùng lượng DNA khác nhau. Sau 50 lần huấn luyện và kiểm thử mười classifier, random forest với 19 đặc trưng được chọn làm RFC19 \cite{benschop2019,alamoudi2022}. Kruijver và cộng sự phát triển decision tree minh bạch trên 766 mixture PROVEDIt và cho thấy hiệu năng phụ thuộc mạnh vào bộ lọc stutter cùng artefact; khoảng accuracy được TAWSEEM tổng hợp cho thiết lập này là 0,77 đến 0,85 \cite{kruijver2021,alamoudi2022}. Ưu điểm của decision tree là tốc độ và khả năng giải thích trực tiếp, nhưng việc tóm tắt hồ sơ thành số ít đặc trưng có thể làm mất cấu trúc peak.


\subsection{Deep learning cho Number of Contributors}

TAWSEEM là một MLP sâu cho NOC được huấn luyện trên PROVEDIt. Tác giả mô tả PROVEDIt có hơn 25.000 multiplex STR profile từ một đến năm người và thử ba kịch bản gồm 780 profile một multiplex, 5.700 profile ba multiplex và 6.000 profile bốn multiplex. Kịch bản lớn nhất sử dụng khoảng 6.000 đặc trưng và đạt 97\% accuracy; tác giả cũng báo cáo 97\% accuracy cho mixture năm contributor \cite{alamoudi2022}. Đây là kết quả nổi bật, nhưng đầu vào phẳng phụ thuộc mạnh vào multiplex và khó chuyển trực tiếp sang kit có số locus khác.

deepNoC giải quyết hạn chế dữ liệu bằng mô phỏng electrophoretic signal. Công trình sinh 100.000 profile gán nhãn trước, huấn luyện mạng sâu cho nhãn từ một đến mười contributor và báo cáo 89\% accuracy. Tác giả cho thấy có thể fine tune với vài trăm profile để đạt accuracy tương tự trong một phòng thí nghiệm cụ thể, đồng thời bổ sung các đầu ra phụ nhằm hỗ trợ giải thích \cite{taylor2024deepnoc}. So với TAWSEEM, deepNoC mở rộng miền nhãn nhưng phụ thuộc chất lượng của simulator; khoảng cách giữa synthetic signal và casework thật vẫn cần được định lượng.


Veldhuis và cộng sự tập trung vào explainability. Thay vì chỉ đưa ra feature importance, họ xây dựng counterfactual explanation thực tế cho một dự đoán NOC, nghĩa là tạo thay đổi có thể xảy ra trong hồ sơ DNA và đủ làm mô hình đổi nhãn \cite{veldhuis2022}. Hướng này đặc biệt quan trọng vì một classifier đạt accuracy cao nhưng không giải thích được yếu tố peak hoặc locus nào dẫn đến quyết định sẽ khó được kiểm chứng trong môi trường pháp y.


\subsection{Deep learning cho nhận dạng contributor và allele calling}

Phan và cộng sự đề xuất pipeline học sâu dự đoán cá thể trong mixture sequencing, sử dụng 27 STR và 94 SNP trên dữ liệu từ sáu cá thể \cite{phan2021}. Bài toán này gần contributor identification hơn NoC, nhưng dữ liệu massively parallel sequencing giàu thông tin sequence hơn capillary electrophoresis. Vì vậy, hiệu năng của công trình không thể so sánh trực tiếp với mô hình xử lý peak height.

DNANet chuyển allele calling thành segmentation trên ma trận electropherogram kích thước năm dye lane nhân 4.096 scan point. Tác giả khảo sát 90 cấu hình tạo bởi ba độ sâu mạng, hai lựa chọn kernel và năm số filter khởi đầu. Phiên bản xuất bản báo cáo F1 bằng 0,971 trên allele do analyst gán nhãn trong case data chưa thấy, 0,982 trên research data và 0,962 khi so với allele donor thật, bằng F1 của annotation analyst trên tiêu chí này \cite{dewit2025}. DNANet không dự đoán NOC, nhưng cho thấy việc giữ cấu trúc tín hiệu thay vì chỉ dùng feature tổng hợp có giá trị thực nghiệm.

Khung đề xuất trong nghiên cứu này nằm sau bước allele calling. Đầu vào là các peak đã được trích xuất và chuẩn hóa, trong khi DNANet có thể trở thành module phía trước để tự động tạo peak đáng tin cậy. Sự kết hợp tiềm năng giữa DNANet và DS-ST hình thành pipeline từ electropherogram thô đến allele, NOC và contributor, nhưng việc tích hợp đầy đủ nằm ngoài phạm vi thực nghiệm hiện tại.

\subsection{Mạng học trên tập hợp và attention phân cấp}

Transformer chuẩn sử dụng scaled dot product attention
\begin{equation}
\operatorname{Attention}(Q,K,V)=
\operatorname{softmax}\left(\frac{QK^\top}{\sqrt{d_k}}\right)V,
\end{equation}
và loại bỏ recurrence trong kiến trúc gốc \cite{vaswani2017}. Khi không sử dụng positional encoding, self attention có tính equivariant theo hoán vị, phù hợp với tập peak. Set Transformer xây dựng Multihead Attention Block, Induced Set Attention Block và Pooling by Multihead Attention để học hàm trên tập \cite{lee2019}. SetNorm bổ sung chuẩn hóa dành riêng cho tập và clean path residual, giúp các mạng set sâu ổn định hơn \cite{zhang2023setnorm}.

Deep Sets tạo biểu diễn bất biến bằng phép cộng
\begin{equation}
\bm{z}=\rho\left(\sum_{i=1}^{n}\phi(\bm{x}_i)\right).
\end{equation}
Trong mô hình của đồ án, PMA thay phép cộng cố định bằng một seed vector học được. Seed này truy vấn toàn bộ tập peak đã mã hóa và tạo biểu diễn hồ sơ 128 chiều. Nhờ không sử dụng positional encoding theo thứ tự cột, phép mã hóa giữ tính equivariant theo hoán vị ở mức token, còn PMA tạo đầu ra invariant theo hoán vị. Đây là điểm khác biệt quan trọng với CNN một chiều và với Transformer chia cột locus trong baseline.

\subsection{Slot attention và optimal transport cho contributor identification}

Slot Attention tạo một tập véc tơ slot và cập nhật chúng qua nhiều vòng attention cạnh tranh \cite{locatello2020slot}. Trong thị giác, mỗi slot thường biểu diễn một vật thể; trong nghiên cứu này, mỗi slot được gắn với một donor tham chiếu. Conditional Slot Attention sử dụng grounded slot dictionary để khởi tạo slot có điều kiện thay vì lấy mẫu hoàn toàn ngẫu nhiên \cite{ma2024cosa}. Tương tự, genotype của từng donor được mã hóa thành slot khởi tạo, giúp donor slot bắt đầu từ allele tham chiếu của chính donor đó.

Guided Slot Attention sử dụng thông tin hướng dẫn để cải thiện phân tách foreground và background trong video \cite{lee2024gsa}. Khi điều chỉnh sang DNA, logits gán peak cho donor đóng vai trò hướng dẫn. MESH liên hệ Slot Attention với optimal transport và tối thiểu hóa entropy của Sinkhorn để phá hòa giữa các slot tương tự \cite{zhang2023mesh}. Điều này phù hợp với trường hợp hai donor chia sẻ nhiều allele. Adaptive Slot Attention giải quyết số slot biến đổi bằng cơ chế chọn slot rời rạc và masked decoder \cite{fan2024adaslot}; trong bối cảnh DNA, gate tồn tại của slot là tín hiệu trực tiếp cho việc donor có mặt hay không.

\subsection{Khoảng trống nghiên cứu}
\label{sec:gap}

Tổng hợp các công trình từ ba thế hệ phương pháp cho thấy lĩnh vực NOC đã đạt nhiều cột mốc đáng kể, song bốn khoảng trống vẫn còn rõ ràng và chưa được giải quyết triệt để trong bất kỳ công trình đơn lẻ nào.

\paragraph{Biểu diễn đầu vào không khai thác cấu trúc peak nội locus.}
PACE, RFC19 và TAWSEEM đều chuyển hồ sơ STR thành véc tơ phẳng chiều cao allele trước khi đưa vào classifier. Phép chiếu này xóa đi cấu trúc bên trong từng locus: hai allele cùng locus có quan hệ sinh học chặt chẽ — stutter xuất hiện theo khoảng cách xác định, cân bằng chiều cao phản ánh tỷ lệ contributor, xếp hạng peak tương quan với số locus — song toàn bộ thông tin này biến mất khi xếp phẳng thành cột liền kề. deepNoC cải thiện bằng cách tiếp nhận tín hiệu điện di từ simulator, nhưng bản thân mô hình phân loại vẫn không khai thác cấu trúc locus-aware ở mức kiến trúc. Một biểu diễn tập hợp phân cấp, trong đó mỗi peak là phần tử của locus và mỗi locus là phần tử của profile, có thể giữ được thông tin cấu trúc này và tận dụng tính bất biến hoán vị tự nhiên của dữ liệu STR.

\paragraph{Rò rỉ dữ liệu do thiếu GroupKFold theo mixture identity.}
Trong PROVEDIt, một mixture vật lý thường xuất hiện ở nhiều injection time và điều kiện thu nhận khác nhau. Nếu phân chia ngẫu nhiên theo hàng, các bản gần-trùng lặp của cùng một mixture sẽ xuất hiện đồng thời ở tập huấn luyện và tập validation, tạo ra ước lượng hiệu năng quá lạc quan so với casework thật. Trong số ba công trình tiêu biểu, không có công trình nào mô tả rõ việc dùng GroupKFold kiểm soát mixture identity; điều này đặt nghi vấn về tính thực tiễn của accuracy được báo cáo khi áp dụng mô hình vào hồ sơ chưa gặp. Nghiên cứu này áp dụng GroupKFold năm phần theo \texttt{group\_id} nhất quán ở mọi mô hình và mọi thí nghiệm, đảm bảo không có mixture identity nào xuất hiện ở cả hai phía của ranh giới phân chia.

\paragraph{Thiếu backbone dùng chung cho NOC và contributor identification.}
NOC và contributor identification là hai nhiệm vụ có quan hệ mật thiết: biết số người đóng góp giúp thu hẹp không gian tìm kiếm contributor, trong khi thông tin genotype donor có thể phản chiếu ngược về số lượng người. Hầu hết các công trình đã công bố tập trung riêng lẻ vào một trong hai nhiệm vụ và không đề xuất kiến trúc nào có thể chia sẻ tham số giữa hai đầu nhiệm vụ. Một backbone Set Transformer với tiền huấn luyện trên contributor identification 45 lớp rồi fine-tune cho NOC có thể khai thác mối liên hệ này một cách tường minh, vừa tăng hiệu quả tham số vừa cung cấp khả năng tái sử dụng kiến trúc.

\paragraph{Khả năng giải thích và kiểm chứng còn hạn chế.}
Veldhuis và cộng sự đã chỉ ra rằng feature importance truyền thống không đủ để giải thích quyết định NOC trong môi trường pháp y; họ đề xuất counterfactual explanation đòi hỏi thay đổi thực tế trong hồ sơ thay vì chỉ đảo chiều một đặc trưng trừu tượng \cite{veldhuis2022}. Barash và cộng sự nhấn mạnh thêm nhu cầu về khả năng kiểm chứng liên ngành và hiểu biết chuyên sâu giữa thống kê pháp y và machine learning \cite{barash2024}. Một kiến trúc attention-based cung cấp attention map theo locus và theo donor một cách tự nhiên, tạo điều kiện cho việc kiểm tra bằng chứng ở cấp peak; tuy nhiên, việc sử dụng attention map như bằng chứng pháp y đòi hỏi validation độc lập chuyên sâu hơn và nằm ngoài phạm vi của nghiên cứu này.

Đề xuất trong nghiên cứu này nhắm trực tiếp vào ba khoảng trống đầu: biểu diễn tập hợp phân cấp khai thác cấu trúc peak nội locus, phân chia dữ liệu GroupKFold nghiêm ngặt ngăn rò rỉ mixture identity, và kiến trúc backbone dùng chung cho cả NOC lẫn contributor identification qua cơ chế tiền huấn luyện.

\section{Mô hình đề xuất}
\label{sec:model}

\subsection{Tập dữ liệu}

\subsubsection{Nguồn dữ liệu PROVEDIt}

PROVEDIt là viết tắt của Project Research Openness for Validation with Empirical Data. Bài báo gốc mô tả một tập dữ liệu quy mô lớn gồm hồ sơ STR nguồn đơn và hỗn hợp để hỗ trợ phát triển, so sánh và validation phương pháp nhận dạng người \cite{provedit2018}. Trang chính thức của Rutgers hiện mô tả hơn 27.000 hỗn hợp DNA pháp y được cung cấp mở \cite{proveditweb}. TAWSEEM cho biết dữ liệu được tạo trong hơn bốn năm với 144 điều kiện phòng thí nghiệm, gồm profile từ một đến năm người và các bản raw, filtered cùng unfiltered \cite{alamoudi2022}.

Nghiên cứu sử dụng ba nhánh dữ liệu được tạo từ các kit GlobalFiler, Identifiler Plus và PowerPlex 16. GlobalFiler gồm 10.124 hồ sơ và 439 cột tính cả định danh nhóm cùng nhãn, Identifiler Plus gồm 3.160 hồ sơ và 281 cột, còn PowerPlex 16 gồm 1.022 hồ sơ và 214 cột. Sự khác biệt số cột phản ánh số locus và allele bin khác nhau giữa các kit, là lý do kiến trúc không nên phụ thuộc một véc tơ đầu vào có chiều cố định.

\begin{table}[H]
\centering
\caption{Thống kê các tập dữ liệu đã xây dựng từ PROVEDIt.}
\label{tab:data_stats}
\small
\begin{tabular}{lrrrrrrr}
\toprule
\textbf{Bộ kit} & \textbf{Hồ sơ} & \textbf{Nhóm} & \textbf{NOC 1} & \textbf{NOC 2} &
\textbf{NOC 3} & \textbf{NOC 4} & \textbf{NOC 5} \\
\midrule
GlobalFiler & 10.124 & 3.418 & 8.119 & 526 & 484 & 527 & 468 \\
Identifiler Plus & 3.160 & 1.084 & 2.231 & 365 & 209 & 146 & 209 \\
PowerPlex 16 & 1.022 & 435 & 794 & 56 & 52 & 60 & 60 \\
\bottomrule
\end{tabular}
\end{table}

Phân phối GlobalFiler có 80,20\% hồ sơ thuộc lớp một contributor, trong khi lớp hiếm nhất là năm contributor chỉ chiếm 4,62\%. Nếu tối ưu accuracy thông thường, mô hình có thể đạt điểm cao bằng cách ưu tiên NOC một nhưng thất bại ở mixture phức tạp. Vì vậy, Macro F1 được sử dụng làm metric lựa chọn chính, Focal Loss giảm trọng số mẫu dễ và oversampling được thực hiện chỉ trong tập huấn luyện.

\subsubsection{Ngăn rò rỉ dữ liệu}

Một mixture vật lý trong PROVEDIt có thể xuất hiện ở nhiều injection time hoặc điều kiện thu nhận. Nếu chia ngẫu nhiên theo hàng, các biến thể gần như trùng nhau của cùng mixture có thể nằm ở cả train và validation, tạo ra đánh giá quá lạc quan. Pipeline xây dựng \texttt{group\_id} bằng cách loại thành phần thời gian tiêm khỏi tên sample, sau đó sử dụng GroupKFold. Với năm fold, mỗi vòng dùng bốn nhóm fold để huấn luyện và một fold để kiểm định; không mixture identity nào xuất hiện ở cả hai phía.

Ablation được chạy bằng GroupKFold ba phần để giảm chi phí, do đó các giá trị ablation chỉ được so với baseline ba phần cùng protocol và không được trộn với kết quả GroupKFold năm phần của benchmark chính.

\subsubsection{Tiền xử lý và phân chia dữ liệu GF29cycles}

Bên cạnh tập GlobalFiler 10.124 hồ sơ dùng cho benchmark NOC (Bài toán 2), nghiên cứu xây dựng thêm một tập dữ liệu quy mô lớn hơn từ nhánh \textit{GlobalFiler 29 cycles} (\textbf{GF29cycles}) của PROVEDIt, phục vụ tiền huấn luyện DS-ST (Bài toán 1) và đánh giá khả năng tổng quát hóa.

\paragraph{Nguồn dữ liệu thô và tạo hỗn hợp.}
Dữ liệu thô là các file CSV chứa profile STR GlobalFiler thu nhận ở 29 chu kỳ PCR, bao gồm cả profile nguồn đơn (NOC\,=\,1) và hỗn hợp hai đến năm người. Để tăng tính đa dạng lớp NOC\,=\,1, các profile một người được trực tiếp đưa vào tập huấn luyện thay vì loại bỏ. Bước \textit{Mix with NOC\,=\,1} kết hợp profile một người từ đủ mọi điều kiện phòng thí nghiệm vào cùng pipeline xử lý với các hỗn hợp, tạo ra tổng cộng \textbf{55.247 hồ sơ} sau bước lọc và chuẩn hóa.

\paragraph{Phân chia train / val / test.}
Tập 55.247 hồ sơ được phân chia thành ba phần theo mixture identity để ngăn rò rỉ dữ liệu, như minh họa trong Hình~\ref{fig:gf29_preprocess}.

\begin{figure}[H]
\centering
\begin{tikzpicture}[
  node distance=0.55cm and 1.2cm,
  gbox/.style={draw, rounded corners=7pt, fill=green!12, align=center,
               font=\small, inner sep=7pt, minimum width=3.0cm},
  arr/.style={-Stealth, thick, gray!60}
]

\node[gbox] (raw) {Raw CSVs\\GF29cycles};
\node[gbox, below=of raw] (mix) {Mix with NOC\,=\,1};
\node[gbox, minimum width=4.2cm, below=of mix] (total)
  {\textbf{55.247 hồ sơ}};

\node[gbox, minimum width=3.5cm,
      below left=1.4cm and 3.2cm of total] (train)
  {\textbf{Train}\\[3pt]
   5.247\enspace NOC\,=\,1\\
   7.266\enspace NOC\,=\,2\\
   10.625\;NOC\,=\,3\\
   17.709\;NOC\,=\,4\\
   14.400\;NOC\,=\,5};

\node[gbox, minimum width=3.0cm,
      below=1.4cm of total] (val)
  {\textbf{Val}\\[3pt]
   3.582 hồ sơ};

\node[gbox, minimum width=3.5cm,
      below right=1.4cm and 3.2cm of total] (test)
  {\textbf{Test}\\[3pt]
   1.366 closed-set\\
   1.366 open-set};

\draw[arr] (raw)   -- (mix);
\draw[arr] (mix)   -- (total);
\draw[arr] (total) -- (train);
\draw[arr] (total) -- (val);
\draw[arr] (total) -- (test);

\end{tikzpicture}
\caption{Sơ đồ tiền xử lý và phân chia dữ liệu GF29cycles. Từ tổng cộng 55.247 hồ sơ sau bước \textit{Mix with NOC\,=\,1}, tập được chia thành ba phần theo mixture identity: huấn luyện, validation và kiểm thử (closed-set và open-set).}
\label{fig:gf29_preprocess}
\end{figure}

\paragraph{Đặc điểm phân phối.}
Khác với GlobalFiler 10.124 hồ sơ (80,2\% là NOC\,=\,1), tập GF29cycles có phân phối lớp đồng đều hơn: lớp nhiều nhất là NOC\,=\,4 chiếm 32,1\% và lớp ít nhất là NOC\,=\,1 chỉ chiếm 9,5\%. Điều này phản ánh sự kết hợp có chủ đích giữa mixture thực và profile một người, tạo điều kiện để DS-ST học nhận dạng contributor ngay cả khi mixture chứa nhiều người tham gia.

\paragraph{Giao thức đánh giá.}
Với kịch bản closed-set, metric chính là Macro F1 trên năm lớp NOC và average precision của contributor binary classifier. Với kịch bản open-set, thêm chỉ số rejection accuracy — tỷ lệ hồ sơ ngoài panel được mô hình từ chối đúng thông qua đầu \texttt{logit\_reject} và ngưỡng sau mô hình. Bảng~\ref{tab:data_stats} tóm tắt thống kê cả hai tập.

\subsection{Biểu diễn đầu vào}

Mỗi hồ sơ được chuyển thành tensor
\begin{equation}
\mathcal{X}\in\mathbb{R}^{B\times L\times P\times F},
\end{equation}
trong đó \(B\) là batch size, \(L=24\) là số locus tối đa của GlobalFiler, \(P=50\) là số peak tối đa tại mỗi locus và \(F=89\) là số đặc trưng peak sau tiền xử lý. Mask
\(\mathcal{M}\in\{0,1\}^{B\times L\times P}\) đánh dấu peak hợp lệ. Khi một kit không có locus tương ứng, toàn bộ locus được mask thay vì đánh lại số locus, nhờ đó định danh sinh học của locus được giữ nhất quán.

Tập đặc trưng bao gồm định danh locus và allele, log chiều cao peak, tỷ lệ chiều cao so với tổng locus, chiều cao tương đối toàn hồ sơ, thứ hạng peak trong locus, số peak tại locus, khoảng cách tới allele lân cận và các thống kê giúp nhận biết stutter. Đối với pipeline DS-ST, phiên bản rút gọn tám chiều được sử dụng:
\begin{equation}
\bm{x}_{lp}=
\left[l,\ a,\ \log(1+h),\ H_b,\ SR,\ r,\ n_{\mathrm{peak}},\ h_{\mathrm{rel}}\right].
\end{equation}
Trong đó \(H_b=h/\sum_{q\in l}h_q\), \(SR=h/h_{a+1}\) khi peak cha tồn tại, \(r\) là rank chuẩn hóa, \(n_{\mathrm{peak}}\) là số peak chuẩn hóa và \(h_{\mathrm{rel}}=h/\max_q h_q\).

\subsection{Hai bài toán và kiến trúc backbone chung}

Nghiên cứu đề xuất hai nhánh của cùng mô hình DS-ST, chia sẻ backbone Set Transformer nhưng phục vụ hai mục tiêu riêng biệt:

Bài toán thứ nhất sử dụng nhánh contributor của DS-ST để xác định tập contributor từ panel tham chiếu $C=45$ donor. Backbone được mở rộng với feasibility filter, set-of-set encoder và bộ giải mã slot có điều kiện genotype theo chuỗi CoSA, GSANet, MESH và AdaSlot.

Bài toán thứ hai sử dụng nhánh NoC của DS-ST để dự đoán số lượng contributor $\mathrm{NOC}\in\{1,\ldots,5\}$. Backbone được fine-tune từ checkpoint contributor, còn bộ giải mã slot được thay bằng đầu NOC nhỏ gọn $128\to64\to5$ và huấn luyện bằng Focal Loss.

Hai nhánh chia sẻ cùng backbone Set Transformer — locus embedding kết hợp Periodic Linear Representation (PLR) \cite{gorishniy2022} cho các đặc trưng số liên tục, hai khối ISAB++ \cite{lee2019}, 32 inducing points, $d_{\mathrm{model}}=128$ — cho phép trọng số chuyển giao từ nhánh contributor sang nhánh NoC qua tiền huấn luyện.

\begin{figure}[H]
\centering
\resizebox{\textwidth}{!}{
\begin{tikzpicture}[node distance=0.6cm and 0.55cm]
  % ── Backbone (hàng chính) ─────────────────────────────────
  \node[block,fill=peakblue!12] (tok)
    {Token tám\\trường\\$(B,N,8)$};
  \node[block,fill=peakblue!12,right=of tok] (enc)
    {Mã hóa token\\LocusEmbed\\$+$ PLR\\$+$ feas\_filter};
  \node[block,fill=locusgreen!14,right=of enc] (isab)
    {ISAB++ $\times$2\\set\_of\_set\\$\mathbf{H}(B,N,128)$};
  % ── Panel tham chiếu (trên) ────────────────────────────────
  \node[block,fill=gray!12,above right=1.6cm and 2.3cm of isab] (panel)
    {Panel kiểu gen\\$C{=}45$ donor};
  % ── Nhánh trên: Bộ giải mã slot ───────────────────────────
  \node[block,fill=headpurple!13,right=2.3cm of isab,yshift=0.45cm] (dec)
    {Bộ giải mã khung\\CoSA $\to$ GSANet\\$\to$ MESH $\to$ AdaSlot};
  \node[smallblock,right=of dec] (sigma)
    {$\sigma(\ell)$\\$P(\mathrm{donor}_c\!\mid\!X)$};
  % ── Nhánh dưới: PMA → reject ──────────────────────────────
  \node[block,fill=profileorange!14,below right=1.1cm and 0.9cm of isab] (pma)
    {PMA\\$\to$ \#\ reject};
  % ── Phi-rerank ────────────────────────────────────────────
  \node[block,fill=locusgreen!12,right=2.8cm of pma,yshift=0.55cm] (phi)
    {phi-rerank\\$+$ đếm RF};
  \node[smallblock,right=of phi] (out)
    {Tập contributor\\$\uparrow$ NOC $k$};
  % ── Arrows ────────────────────────────────────────────────
  \draw[line] (tok)--(enc);
  \draw[line] (enc)--(isab);
  \draw[line] (isab.east)       to[out=10,in=180]  (dec.west);
  \draw[line] (isab.south east) to[out=-25,in=180] (pma.west);
  \draw[line] (panel.south)     to[out=-90,in=90]  (dec.north);
  \draw[line] (dec)--(sigma);
  \draw[line] (sigma.east)      to[out=0,in=90]    (phi.north);
  \draw[line] (pma.east)        to[out=0,in=-90]   (phi.south);
  \draw[line] (phi)--(out);
\end{tikzpicture}}
\caption{\textbf{Bài toán 1 — Luồng tổng quát DS-ST.} Token tám trường $(B,N,8)$ qua LocusEmbed + PLR + feas\_filter, rồi hai khối ISAB++ với set\_of\_set cho biểu diễn $\mathbf{H}(B,N,128)$. Nhánh dưới đến PMA (mixture reject); nhánh trên kết hợp genotype từ panel $C{=}45$ donor qua bộ giải mã slot (CoSA$\to$GSANet$\to$MESH$\to$AdaSlot) trả về $\sigma(\ell_c)$ cho mỗi donor. Hai nhánh hội tụ tại phi-rerank + Random Forest để suy ra tập contributor và NOC cuối.}
\label{fig:dsst_overview}
\end{figure}

Hình~\ref{fig:dsst_overview} thể hiện luồng DS-ST. Backbone xử lý toàn bộ peak của hỗn hợp như một tập bằng phẳng, không có cấu trúc phân cấp locus cứng nhắc trong backbone. Khi triển khai liên kết, $k=\hat{y}^{\mathrm{noc}}$ từ Bài toán 2 được dùng để chọn top $k$ donor; trong đánh giá riêng biệt, oracle top true $k$ được báo cáo để tách chất lượng xếp hạng khỏi lỗi đếm. Chi tiết backbone encoder DS-ST được trình bày trong Hình~\ref{fig:dsst_encoder}.

\begin{figure}[H]
\centering
\includegraphics[width=0.88\linewidth]{figures/fig3_encoder.png}
\caption{\textbf{Chi tiết backbone encoder DS-ST.} Peak được chia thành nhóm private (chỉ của một donor) và shared (chia sẻ giữa nhiều donor) theo feasibility filter (\texttt{set\_of\_set}). Hai khối ISAB++ xử lý toàn bộ tập peak; cardinality estimate $\hat{n}$ từ PMA định tuyến giữa nhánh xử lý hỗn hợp và nhánh reject (hỗn hợp không có trong panel). Trong nhánh NoC, set\_of\_set và reject branch không được sử dụng; mô hình dùng backbone ISAB++ với token tám chiều.}
\label{fig:dsst_encoder}
\end{figure}

\subsection{Bài toán 2 — DS-ST: Dự đoán Number of Contributors}

Nhánh NoC của DS-ST giải quyết Bài toán 2. Toàn bộ peak của hồ sơ được ghép thành một chuỗi token thống nhất và xử lý bằng backbone Set Transformer dùng chung, nhưng không đi qua bộ giải mã slot. Backbone được fine-tune từ checkpoint đã tiền huấn luyện trên nhiệm vụ contributor identification với 45 lớp donor. Hình~\ref{fig:settransv4} mô tả luồng chi tiết của nhánh này.

\begin{figure}[H]
\centering
\resizebox{\textwidth}{!}{
\begin{tikzpicture}[node distance=0.55cm and 0.55cm]
  \node[block,fill=peakblue!12] (input)
    {Token tám chiều\\$(B,N,8)$};
  \node[block,fill=peakblue!12,right=of input] (embed)
    {LocusEmbed\\$+$ PeriodicPLR\\$\to$ InputProj\\$(B,N,128)$};
  \node[block,fill=locusgreen!14,right=of embed] (mab0)
    {mab0\\Sigmoid-\\MABpp};
  \node[block,fill=locusgreen!12,right=of mab0] (mab1)
    {mab1\\MABpp\\softmax};
  \node[block,fill=profileorange!14,right=of mab1] (pma)
    {PMA\\$k{=}1$\\$\bm{z}_{\mathrm{mix}}\ (B,128)$};
  \node[block,fill=headpurple!13,right=of pma] (head)
    {NOC Head\\$128\!\to\!64\!\to\!5$\\Focal Loss};
  \node[smallblock,right=of head] (out)
    {$P(y\!=\!1..5)$};
  \node[groupbox,fit=(mab0)(mab1)] (isab) {};
  \node[above=2pt of isab,font=\tiny\bfseries] {ISABpp $\times$ 2,\ SetNorm,\ clean-path};
  \draw[line] (input)--(embed);
  \draw[line] (embed)--(mab0);
  \draw[line] (mab0)--(mab1);
  \draw[line] (mab1.east)--(pma.west);
  \draw[line] (pma)--(head);
  \draw[line] (head)--(out);
\end{tikzpicture}}
\caption{\textbf{Bài toán 2 — Luồng DS-ST cho NOC.} Token tám chiều mỗi peak được mã hóa song song qua LocusEmbed (16 chiều) và PeriodicPLR (56 chiều) rồi chiếu tuyến tính vào $d_{\mathrm{model}}=128$ (chi tiết embedding xem Hình~\ref{fig:dsst_embed}). Hai khối ISAB++ với SigmoidMABpp tại mab0 học quan hệ giữa mọi peak qua 32 inducing points (chi tiết khối ISAB++ xem Hình~\ref{fig:dsst_isab}). PMA ($k{=}1$) tổng hợp thành $\bm{z}_{\mathrm{mix}}\in\mathbb{R}^{128}$ và đầu NOC chiếu xuống năm lớp bằng Focal Loss.}
\label{fig:settransv4}
\end{figure}

Mỗi peak được biểu diễn bằng token tám chiều
\begin{equation}
\bm{x}_p = \bigl[l,\;a,\;\log(1+h),\;H_b,\;SR,\;r,\;n_l/12,\;h/h_{\max}\bigr],
\end{equation}
trong đó $l$ là chỉ số locus ($0\ldots23$), $a$ là giá trị allele, $h$ là chiều cao peak RFU, $H_b=h/\sum_{q\in l}h_q$ là tỷ lệ chiều cao trong locus, $SR=h/(h_{a+1}+\varepsilon)$ là tỷ số với chiều cao peak stutter cha tại allele $a+1$ (bằng không nếu peak cha không tồn tại), $r=\mathrm{rank}_{\downarrow}/(n_l-1)$ là thứ hạng giảm dần chuẩn hóa trong locus (bằng không cho peak cao nhất), $n_l/12$ là số peak chuẩn hóa và $h/h_{\max}$ là chiều cao tương đối so với peak cao nhất toàn hồ sơ. Bảy trường liên tục được chuẩn hóa z-score theo thống kê tập huấn luyện trước khi đưa vào PeriodicPLR.

Chỉ số locus $l$ được mã hóa bằng $\texttt{Embedding}(25,16,\text{padding\_idx}=24)$ cho ra 16 chiều. Bảy trường liên tục sau khi chuẩn hóa z-score được đưa qua Periodic Linear Representation (PLR) \cite{gorishniy2022} với ($n_{\mathrm{freq}}=8$, $d_{\mathrm{num}}=8$, $\sigma_{\mathrm{init}}=0{,}3$): mỗi đặc trưng $x_j$ sinh $[\sin(2\pi c_{jk}x_j), \cos(2\pi c_{jk}x_j)]_{k=1}^{n_{\mathrm{freq}}}$, chiếu qua Linear$(2n_{\mathrm{freq}}\to d_{\mathrm{num}})$ rồi ReLU, tạo 8 chiều riêng; bảy đặc trưng ghép thành $7\times8=56$ chiều. Tần số $c_{jk}$ được học từ khởi tạo $\mathcal{N}(0,\sigma^2)$ và đặc trưng hóa chu kỳ tuần hoàn tự nhiên của allele, chiều cao và rank. Hai nhánh embedding được ghép thành $16+56=72$ chiều rồi chiếu vào $d_{\mathrm{model}}=128$ qua $W_{\mathrm{in}}\in\mathbb{R}^{128\times72}$; token của các peak được mask được triệt về không trước khi vào encoder.

\begin{figure}[H]
\centering
\includegraphics[width=0.85\linewidth]{figures/fig2_embedding.png}
\caption{\textbf{Chi tiết mã hóa token dùng chung cho hai nhánh DS-ST.} Chỉ số locus sinh embedding 16 chiều; bảy trường số liên tục qua PeriodicPLR sinh 56 chiều; hai nhánh ghép thành 72 chiều rồi chiếu tuyến tính vào $d_{\mathrm{model}}=128$. Feasibility filter chỉ được sử dụng trong nhánh contributor.}
\label{fig:dsst_embed}
\end{figure}

Mỗi khối ISAB++ gồm hai bước. Bước mab0 dùng SigmoidMABpp, trong đó $M=32$ inducing point thu nhận tín hiệu peak qua attention không cạnh tranh
\begin{equation}
A_{mc}=\sigma\!\left(\frac{Q_mK_c^\top}{\sqrt{d_h}}+b_{\mathrm{nc}}\right),
\qquad b_{\mathrm{nc}}=-\log n_{\mathrm{valid}},
\end{equation}
với $\sigma$ là hàm sigmoid và $b_{\mathrm{nc}}$ là bias tự động điều chỉnh theo số peak hợp lệ trong mẫu. Khác với softmax, các gate này độc lập nhau, cho phép peak thiểu số được chấp nhận vào inducing point mà không bị triệt tiêu bởi peak chiếm ưu thế. Bước mab1 dùng MABpp softmax chuẩn để mỗi peak truy vấn ngược biểu diễn inducing. SetNorm chuẩn hóa trên toàn tập token và clean-path residual ổn định gradient khi xếp chồng nhiều khối. \begin{figure}[H]
\centering
\includegraphics[width=0.82\linewidth]{figures/fig4_isab_block.png}
\caption{\textbf{Chi tiết một khối ISAB++ dùng chung cho hai nhánh DS-ST.} mab0 dùng SigmoidMABpp: $M=32$ inducing points thu nhận tín hiệu peak qua sigmoid attention không cạnh tranh. mab1 dùng MABpp softmax chuẩn để peak truy vấn ngược biểu diễn inducing. SetNorm chuẩn hóa trên toàn tập token và clean-path residual ổn định gradient khi xếp chồng nhiều khối.}
\label{fig:dsst_isab}
\end{figure}

Sau hai ISAB++, PMA với một seed vector học được tạo $\bm{z}_{\mathrm{mix}}\in\mathbb{R}^{128}$. Đầu NOC áp dụng hai lớp tuyến tính với dropout và ReLU:
\begin{equation}
\hat{\bm{p}} = \operatorname{softmax}\!\bigl(W_2\operatorname{ReLU}(W_1\bm{z}_{\mathrm{mix}})\bigr),
\qquad W_1\in\mathbb{R}^{64\times128},\;W_2\in\mathbb{R}^{5\times64}.
\end{equation}
Quá trình fine tune gồm hai giai đoạn với AdamW ($\text{lr}=3\times10^{-4}$, $\lambda=10^{-4}$) và gradient clipping $\|\nabla\|_{\max}=1{,}0$. Giai đoạn một đóng băng toàn bộ tham số backbone và cập nhật chỉ đầu NOC với learning rate $9\times10^{-4}$ trong năm epoch, ổn định giá trị khởi tạo ngẫu nhiên trước khi mở backbone. Giai đoạn hai mở toàn bộ tham số, áp dụng CosineAnnealingLR suốt 60 epoch với $\eta_{\min}=10^{-6}$ và dừng sớm theo Macro F1 validation với patience bằng mười hai. Batch size bằng 64; trước mỗi fold, tập huấn luyện được oversample để cân bằng năm lớp NOC. Focal Loss $\mathcal{L}_{\mathrm{focal}} = -\alpha_y(1-p_y)^{\gamma}\log p_y$ với $\gamma=2$ và trọng số lớp $\alpha_y = N_{\mathrm{total}}/(5\cdot N_y)$ tỉ lệ nghịch tần suất từng lớp.

\subsection{Peak Embedder}

Peak Embedder gồm hai lớp tuyến tính, GELU và LayerNorm. Với peak \(\bm{x}_{lp}\in\mathbb{R}^{89}\),
\begin{align}
\bm{u}^{(1)}_{lp}&=\operatorname{LN}\left(\operatorname{GELU}(W_1\bm{x}_{lp}+\bm{b}_1)\right),\\
\bm{u}^{(2)}_{lp}&=\operatorname{LN}\left(\operatorname{GELU}(W_2\bm{u}^{(1)}_{lp}+\bm{b}_2)\right),
\end{align}
trong đó \(W_1\in\mathbb{R}^{96\times89}\) và \(W_2\in\mathbb{R}^{96\times96}\). Đầu ra có kích thước \(B\times24\times50\times96\). GELU được chọn vì cung cấp hàm kích hoạt trơn và đã được sử dụng rộng rãi trong Transformer \cite{hendrycks2016}.

Trong nhánh DS-ST, bảy đặc trưng số của token tám chiều còn được mã hóa bằng periodic numerical embedding. Với đặc trưng số \(x_j\), mỗi tần số học được \(c_{jk}\) tạo cặp
\begin{equation}
\left[\sin(2\pi c_{jk}x_j),\cos(2\pi c_{jk}x_j)\right],
\end{equation}
sau đó được chiếu sang embedding riêng cho từng feature. Gorishniy và cộng sự cho thấy embedding số theo từng feature có thể cải thiện đáng kể mô hình tabular so với đưa scalar trực tiếp vào backbone \cite{gorishniy2022}.

\subsection{Locus Encoder với Stutter Bias Attention}

Mỗi locus chứa tối đa 50 peak. Self attention chuẩn tính score giữa peak \(i\) và \(j\) từ tích vô hướng query và key. Kiến trúc phân cấp bổ sung bias học được từ khoảng cách allele:
\begin{equation}
s_{ij}^{(h)}=
\frac{\left(\bm{q}_i^{(h)}\right)^\top\bm{k}_j^{(h)}}{\sqrt{d_h}}
g_h\left(|a_i-a_j|\right),
\end{equation}
trong đó \(g_h\) là MLP riêng theo head hoặc MLP cho véc tơ nhiều head. Khoảng cách một repeat unit là tín hiệu quan trọng vì stutter thường xuất hiện gần allele cha. Attention weight được tính sau khi mask padding:
\begin{equation}
\alpha_{ij}^{(h)}=
\frac{\exp(s_{ij}^{(h)})}
{\sum_{j'\in\mathcal{V}_l}\exp(s_{ij'}^{(h)})}.
\end{equation}

Mỗi Locus Encoder block gồm pre LayerNorm, bốn head attention, residual, LayerNorm thứ hai và FFN \(96\rightarrow384\rightarrow96\) với GELU. Hai block được xếp chồng. Đầu ra vẫn có kích thước \(B\times24\times50\times96\), do đó mỗi peak có biểu diễn đã nhận ngữ cảnh từ các peak cùng locus.

\subsection{Deep Sets Pool tại locus}

Với các peak hợp lệ \(\mathcal{V}_l\) của locus \(l\), một MLP \(\phi\) được áp dụng cho từng peak. Ba thống kê được ghép:
\begin{equation}
\bm{s}_l=
\left[
\sum_{p\in\mathcal{V}_l}\phi(\bm{u}_{lp});
\max_{p\in\mathcal{V}_l}\phi(\bm{u}_{lp});
\frac{|\mathcal{V}_l|}{P}
\right].
\end{equation}
Locus token được tính bằng
\begin{equation}
\bm{z}_l=\rho(\bm{s}_l)\in\mathbb{R}^{96}.
\end{equation}
Phép tổng và cực đại không phụ thuộc thứ tự peak. Thành phần count cung cấp thông tin trực tiếp về độ phức tạp locus, nhưng được chuẩn hóa và kết hợp với biểu diễn học được thay vì dùng như quy tắc cứng.

\subsection{Cross Locus Transformer và Profile Pool}

Mỗi locus token được cộng locus embedding và dye embedding. Hai Transformer block cho phép locus tương tác toàn cục. Nếu một locus có nhiều peak do stutter trong khi các locus khác có cấu trúc nguồn đơn rõ ràng, attention toàn hồ sơ có thể giảm ảnh hưởng của locus bất thường. Đầu ra
\(\bm{Z}\in\mathbb{R}^{B\times24\times96}\) được tổng hợp:
\begin{equation}
\bm{z}_{\mathrm{profile}}=
W_p\left[
\sum_l\bm{z}_l;
\frac{1}{|\mathcal{L}|}\sum_l\bm{z}_l;
\max_l\bm{z}_l
\right]+\bm{b}_p,
\end{equation}
trong đó \(\bm{z}_{\mathrm{profile}}\in\mathbb{R}^{B\times96}\). Tổng phản ánh khối lượng profile, mean giúp ổn định theo số locus hợp lệ và max giữ locus có bằng chứng mạnh.

\subsection{Đầu dự đoán Number of Contributors}

Shared MLP biến đổi profile embedding:
\begin{equation}
\bm{h}=\operatorname{Dropout}\left(
\operatorname{GELU}(W_s\bm{z}_{\mathrm{profile}}+\bm{b}_s)
\right),\qquad \bm{h}\in\mathbb{R}^{96}.
\end{equation}
Đầu phân loại tạo logits \(\bm{o}^{\mathrm{cls}}\in\mathbb{R}^{5}\) và
\(\bm{p}^{\mathrm{cls}}=\operatorname{softmax}(\bm{o}^{\mathrm{cls}})\).

Đầu CORN tạo bốn logits ứng với các điều kiện \(Y>1,\ldots,Y>4\). Xác suất vượt ngưỡng được tính:
\begin{equation}
P(Y>k\mid\bm{x})=
\prod_{j=1}^{k}
P(Y>j\mid Y>j-1,\bm{x}).
\end{equation}
Từ các xác suất tích lũy, mô hình suy ra phân bố \(\bm{p}^{\mathrm{corn}}\) trên năm lớp. Việc dùng conditional training set đảm bảo rank consistency mà không buộc các output task chia sẻ trọng số như CORAL \cite{shi2023corn}.

Đầu hồi quy tạo scalar \(\hat{n}\in\mathbb{R}\), sau đó chuyển thành phân bố mềm quanh các tâm \(1,\ldots,5\):
\begin{equation}
p_k^{\mathrm{reg}}=
\frac{\exp\left(-(\hat{n}-k)^2/(2\sigma^2)\right)}
{\sum_{j=1}^{5}\exp\left(-(\hat{n}-j)^2/(2\sigma^2)\right)}.
\end{equation}
Phân bố cuối là trung bình:
\begin{equation}
\bm{p}^{\mathrm{noc}}=
\frac{\bm{p}^{\mathrm{cls}}+\bm{p}^{\mathrm{corn}}+\bm{p}^{\mathrm{reg}}}{3},
\qquad
\hat{y}^{\mathrm{noc}}=\arg\max_k p_k^{\mathrm{noc}}.
\end{equation}

Một đầu phụ dự đoán số allele tại mỗi locus với tensor \(B\times24\times20\). Loss phụ ép locus token giữ thông tin đếm cục bộ. Đầu mixture proportion tạo \(\hat{\bm{\pi}}\in\mathbb{R}^{5}\) qua MLP \(96\rightarrow48\rightarrow5\) và softmax. Vì không có ánh xạ thứ tự donor cố định trong bài toán chỉ đếm, tỷ lệ được sắp giảm dần trước khi tính loss.

\subsection{Bộ giải mã contributor identification (DS-ST — Bài toán 1)}

\begin{figure}[H]
\centering
\includegraphics[width=0.88\linewidth]{figures/fig5_decoder.png}
\caption{\textbf{Bộ giải mã DS-ST.} Biểu diễn peak $\mathbf{H}(B,N,128)$ từ backbone kết hợp với anchor genotype của $C=45$ donor để khởi tạo slot theo nguyên lý Conditional Slot Attention (CoSA \cite{ma2024cosa}). Slot được tinh chỉnh theo tín hiệu attribution lấy cảm hứng từ Guided Slot Attention (GSANet \cite{lee2024gsa}), phân phối peak về donor qua Sinkhorn optimal transport theo MESH \cite{zhang2023mesh} và gate tồn tại Gumbel-Sigmoid theo Adaptive Slot Attention (AdaSlot \cite{fan2024adaslot}). Đầu ra là xác suất $\sigma(\ell_c)$ cho mỗi donor.}
\label{fig:dsst_decoder}
\end{figure}

Kiến trúc bộ giải mã được mô tả trong Hình~\ref{fig:dsst_decoder}. Với panel \(C=45\) donor, genotype tham chiếu của donor \(c\) được mã hóa bằng cùng backbone và pooling để tạo anchor \(\bm{g}_c\in\mathbb{R}^{128}\). Các anchor được broadcast thành
\(\bm{S}^{(0)}\in\mathbb{R}^{B\times C\times128}\). Đây là phiên bản điều chỉnh của Conditional Slot Attention, trong đó grounded slot dictionary được thay bằng genotype dictionary \cite{ma2024cosa}.

Đầu attribution tạo logits
\(\bm{A}^{\mathrm{attr}}\in\mathbb{R}^{B\times N\times(C+1)}\), với cột cuối là background. Softmax attribution được dùng để tổng hợp peak cho từng donor:
\begin{equation}
\bm{r}_c=\sum_{i=1}^{N}
\operatorname{softmax}(\bm{A}^{\mathrm{attr}})_{ic}\bm{H}_i.
\end{equation}
Một gate học được quyết định mức cập nhật \(\bm{S}^{(0)}_c+\sigma(w^\top\bm{r}_c)\bm{r}_c\). Cơ chế này điều chỉnh ý tưởng Guided Slot Attention \cite{lee2024gsa} bằng tín hiệu tương thích giữa peak quan sát và genotype tham chiếu của donor — thay vì tín hiệu foreground/background như trong bài toán gốc về video segmentation.

Trong mỗi vòng MESH, affinity giữa slot và peak là
\begin{equation}
\bm{C}=\frac{Q(\bm{S})K(\bm{H})^\top}{\sqrt{d}},
\qquad \bm{C}\in\mathbb{R}^{B\times C\times N}.
\end{equation}
Sinkhorn chuẩn hóa luân phiên theo hàng và cột để tạo ma trận vận chuyển gần doubly stochastic. MESH được đề xuất để kết hợp tốc độ của regularized optimal transport với khả năng phá hòa của unregularized transport \cite{zhang2023mesh}. Nghiên cứu sử dụng ba vòng cập nhật slot, \(\epsilon=0,05\) và năm bước Sinkhorn mỗi vòng. Update được đưa qua GRU và FFN.

Adaptive Slot Attention (AdaSlot) \cite{fan2024adaslot} tạo existence logit \(e_c\) và gate
\begin{equation}
g_c=\sigma\left((e_c+\eta_c)/\tau\right),
\end{equation}
trong đó \(\eta_c\) là logistic noise khi huấn luyện và \(\tau=1\). Contributor logit kết hợp nội dung slot và tồn tại:
\begin{equation}
\ell_c=\ell_c^{\mathrm{content}}+e_c.
\end{equation}
Đầu ra xác suất là \(\sigma(\ell_c)\). Khi suy luận liên kết, mô hình chọn \(k=\hat{y}^{\mathrm{noc}}\) donor có score cao nhất.

\subsection{Phi reranking (DS-ST — Bài toán 1)}

\begin{figure}[H]
\centering
\includegraphics[width=0.78\linewidth]{figures/fig7_decode_posthoc.png}
\caption{\textbf{Giải mã hậu kỳ DS-ST.} Sau khi bộ giải mã slot trả về xác suất donor, phi-rerank chạy EM trên peak height để ước lượng mixture proportion $\pi_c$, sau đó kết hợp logarithmic opinion pool $R_c = z(\ell_c) + \alpha z(\log(\pi_c+\varepsilon))$. Random Forest độc lập đếm contributor từ allele pattern. Hai nguồn tín hiệu bổ sung nhau cho kết quả tập contributor cuối cùng.}
\label{fig:dsst_posthoc}
\end{figure}

Quy trình giải mã hậu kỳ được mô tả trong Hình~\ref{fig:dsst_posthoc}. Một tín hiệu độc lập về mixture proportion được ước lượng bằng EM từ peak height và genotype compatibility. Nếu peak \(i\) có tập carrier \(\mathcal{C}_i\), trách nhiệm của donor \(c\) được cập nhật:
\begin{equation}
r_{ic}=
\frac{\mathbf{1}[c\in\mathcal{C}_i]\pi_c}
{\pi_{\mathrm{bg}}\exp(-2)+\sum_{j\in\mathcal{C}_i}\pi_j}.
\end{equation}
Mixture proportion mới tỷ lệ với tổng chiều cao đã phân trách nhiệm. Sau mười vòng, score contributor được rerank bằng logarithmic opinion pool:
\begin{equation}
R_c=z(\ell_c)+\alpha z(\log(\pi_c+\varepsilon)).
\end{equation}
Hệ số \(\alpha\) được chọn trên validation, không dùng test. Trên checkpoint contributor seed 42, reranking sơ bộ làm oracle exact match của NOC năm tăng từ 0,831 lên 0,901; vì đây là một checkpoint, kết quả được xem là bằng chứng ban đầu chứ chưa phải kết luận đa seed.

\subsection{Hàm mất mát}

Loss NOC là tổ hợp:
\begin{equation}
\mathcal{L}_{\mathrm{noc}}=
\lambda_{\mathrm{cls}}\mathcal{L}_{\mathrm{focal}}+
\lambda_{\mathrm{corn}}\mathcal{L}_{\mathrm{corn}}+
\lambda_{\mathrm{reg}}\lVert \hat{n}-n\rVert_1+
\lambda_{\mathrm{mix}}\mathcal{L}_{\mathrm{mix}}+
\lambda_{\mathrm{aux}}\mathcal{L}_{\mathrm{allele}}.
\end{equation}
Focal Loss có dạng
\begin{equation}
\mathcal{L}_{\mathrm{focal}}=
-\alpha_y(1-p_y)^\gamma\log p_y,
\end{equation}
với \(\gamma=2\). Lin và cộng sự đề xuất Focal Loss để giảm ảnh hưởng của các mẫu dễ trong bối cảnh mất cân bằng mạnh \cite{lin2017}.

Loss contributor sử dụng Asymmetric Loss cho véc tơ donor, binary cross entropy cho open set reject, cross entropy cho cardinality và hai loss phụ attribution cùng phi:
\begin{equation}
\mathcal{L}_{\mathrm{id}}=
\mathcal{L}_{\mathrm{ASL}}+
0,5\mathcal{L}_{\mathrm{reject}}+
0,3\mathcal{L}_{\mathrm{card}}+
\mathcal{L}_{\mathrm{attr}}+
\mathcal{L}_{\phi}.
\end{equation}
Attribution target của peak được xây dựng từ tích mixture proportion và copy number genotype, sau đó chuẩn hóa trên các donor có khả năng mang allele.

\subsection{Độ đo đánh giá}

Với lớp \(k\), precision và recall được tính:
\begin{equation}
\operatorname{Precision}_k=\frac{TP_k}{TP_k+FP_k},
\qquad
\operatorname{Recall}_k=\frac{TP_k}{TP_k+FN_k}.
\end{equation}
F1 của lớp là
\begin{equation}
F1_k=\frac{2\operatorname{Precision}_k\operatorname{Recall}_k}
{\operatorname{Precision}_k+\operatorname{Recall}_k}.
\end{equation}
Macro F1 là trung bình trên năm lớp, do đó mỗi NOC có trọng số bằng nhau. Micro F1 gộp quyết định của toàn bộ mẫu và trong phân loại đơn nhãn bằng accuracy. Do phân phối NOC một chiếm hơn 80\% ở GlobalFiler, Macro F1 là chỉ số chính.

Đối với contributor identification, exact match yêu cầu toàn bộ véc tơ dự đoán bằng nhãn. Oracle top \(k\) sử dụng NOC thật để đo chất lượng ranking, còn deployed top \(k\) sử dụng NOC dự đoán. Recall at \(k\) đo tỷ lệ contributor thật nằm trong \(k\) vị trí đầu. Reject AUROC đo khả năng phân biệt mixture chỉ gồm donor trong panel với mixture có contributor ngoài panel.

\section{Thử nghiệm và đánh giá}

\subsection{Thiết lập thực nghiệm}

Toàn bộ mô hình được đánh giá trên tập GlobalFiler theo giao thức GroupKFold năm phần dựa trên \texttt{group\_id} của mixture. Việc phân chia theo nhóm thay vì theo hàng là bắt buộc để ngăn rò rỉ dữ liệu: một mixture vật lý có thể xuất hiện ở nhiều injection time, và nếu chia ngẫu nhiên, các bản gần-trùng lặp của cùng mixture sẽ nằm ở cả hai phía của ranh giới phân chia. Trước mỗi fold, tập huấn luyện được cân bằng lớp bằng kỹ thuật kết hợp: oversampling thiểu số và giới hạn đa số ở mức tối đa ba lần kích thước lớp nhỏ nhất, đảm bảo mỗi lớp NOC có đại diện tương đương trong quá trình huấn luyện mà không loại bỏ mẫu từ tập validation. Metric đánh giá chính là Macro F1, tức giá trị trung bình không trọng số của F1 trên năm lớp NOC, được bổ sung bởi Micro F1 phản ánh accuracy tổng thể theo số mẫu.

Để đánh giá khách quan đóng góp của biểu diễn tập hợp phân cấp, ba mô hình nền được triển khai với cùng giao thức phân chia dữ liệu. Mỗi mô hình nền đại diện cho một lớp biểu diễn và một triết lý kiến trúc khác nhau. Hình~\ref{fig:xgb_arch}, \ref{fig:cnn_arch} và \ref{fig:ft_arch} minh họa kiến trúc của từng mô hình nền.

\subsubsection{Cấu hình tính toán}

Toàn bộ quá trình huấn luyện và đánh giá mô hình được thực hiện trên nền tảng Kaggle với một GPU NVIDIA Tesla T4 (16~GB VRAM). Môi trường này đảm nhiệm cả hai pha tính toán chính. Thứ nhất là tiền huấn luyện backbone DS-ST trên nhiệm vụ contributor identification 45 lớp với tập GF29cycles gồm 50.000 hồ sơ bán-tổng hợp — checkpoint sau đó được dùng để khởi tạo nhánh NOC — cùng các mô hình DS-ST và Hybrid của nhánh nhận dạng contributor (Bài toán 1). Thứ hai là toàn bộ benchmark Number of Contributors (Bài toán 2), gồm bốn mô hình nền XGBoost, CNN một chiều, Forensic Transformer, TabPFN và bước fine-tune hai giai đoạn của nhánh NOC, được huấn luyện và đánh giá theo GroupKFold năm phần trên 10.124 hồ sơ GlobalFiler. Các mô hình học sâu (CNN, Forensic Transformer, DS-ST) được hiện thực bằng PyTorch với CUDA; XGBoost và TabPFN dùng thư viện chính thức tương ứng, trong đó TabPFN nạp checkpoint tiền huấn luyện công khai và chỉ thực hiện in-context learning tại thời điểm suy luận.

Một máy tính xách tay cá nhân trang bị GPU NVIDIA GeForce RTX 3050 Laptop (4~GB VRAM) được sử dụng cho các tác vụ nhẹ hơn: tiền xử lý dữ liệu, trích xuất và làm giàu đặc trưng peak, cùng quá trình phát triển và gỡ lỗi mô hình trước khi chạy huấn luyện quy mô đầy đủ trên Kaggle T4.

\subsubsection{Mô hình nền XGBoost}

XGBoost thực thi gradient boosting trên cây quyết định, trong đó một tập hợp $T$ cây được xây dựng tuần tự: mỗi cây mới học cách bù đắp sai số dự đoán tổng hợp của $T-1$ cây trước thông qua gradient của hàm mất mát \cite{chen2016}. Đầu vào là véc tơ phẳng chiều cao allele của 24 locus GlobalFiler — biểu diễn đơn giản nhất và trực tiếp nhất của hồ sơ STR dưới dạng tabular data. Tập hợp gồm 300 cây với độ sâu tối đa sáu cấp mỗi cây; learning rate nhỏ kết hợp với số cây lớn thực hiện shrinkage để kiểm soát overfitting. Cấu hình này phù hợp với khuyến nghị thiết lập gradient boosting cho dữ liệu bảng nhiều chiều \cite{chen2016}.

XGBoost được lựa chọn làm baseline vì tính cạnh tranh được ghi nhận rộng rãi trên tabular data và vì phương pháp này đại diện cho nhóm công trình trước đây sử dụng đặc trưng phẳng như PACE \cite{marciano2017} và RFC19 \cite{benschop2019}. Kết quả XGBoost đặt ra ngưỡng cần vượt qua cho các kiến trúc phức tạp hơn.

\begin{figure}[H]
\centering
\begin{tikzpicture}[
  box/.style={draw, rounded corners=4pt, minimum height=1.0cm,
              align=center, font=\small, fill=blue!7},
  arr/.style={-Stealth, thick, gray}
]
\node[box, minimum width=2.4cm, fill=orange!15] (inp)
  {Véc tơ allele phẳng ($d$ chiều)};
\node[box, minimum width=1.3cm, right=0.8cm of inp] (t1) {Cây $1$};
\node[box, minimum width=1.3cm, right=0.5cm of t1]  (t2) {Cây $2$};
\node[right=0.3cm of t2, font=\large] (dots) {$\cdots$};
\node[box, minimum width=1.3cm, right=0.3cm of dots] (tn) {Cây $T$};
\node[box, minimum width=2.2cm, right=0.8cm of tn, fill=teal!12] (agg)
  {Tổng hợp có trọng số};
\node[box, minimum width=1.8cm, right=0.7cm of agg, fill=red!10] (out)
  {Softmax $\to$ NOC};
\draw[arr] (inp) -- (t1);
\draw[arr] (t1)  -- (t2);
\draw[arr] (t2)  -- (dots);
\draw[arr] (dots) -- (tn);
\draw[arr] (tn)  -- (agg);
\draw[arr] (agg) -- (out);
\node[below=0.35cm of t2, font=\scriptsize, text=gray]
  {mỗi cây bù đắp sai số của tập hợp trước};
\end{tikzpicture}
\caption{Sơ đồ mô hình nền XGBoost. Tập hợp $T=300$ cây quyết định được xây dựng tuần tự theo gradient của hàm mất mát; mỗi cây học trên phần dư của tập hợp trước.}
\label{fig:xgb_arch}
\end{figure}

\subsubsection{Mô hình nền CNN một chiều}

Mạng tích chập một chiều xử lý véc tơ phẳng chiều cao allele như một chuỗi tín hiệu và học đặc trưng phân cấp thông qua ba khối tích chập xếp chồng (Hình~\ref{fig:cnn_arch}). Khối đầu tiên áp dụng 64 bộ lọc kích thước 3 để trích xuất pattern cục bộ giữa các allele liền kề, theo sau bởi chuẩn hóa batch, kích hoạt ReLU và max pooling để giảm độ phân giải. Khối thứ hai và thứ ba lặp lại cấu trúc tương tự với 128 và 256 kênh, cho phép mạng tổng hợp đặc trưng ở vùng nhận ngày càng rộng hơn. Khối cuối thay max pooling bằng adaptive average pooling để tạo biểu diễn cố định 2048 chiều bất kể chiều dài chuỗi đầu vào. Bộ phân loại cuối là hai lớp fully connected với kích hoạt ReLU và dropout ngăn overfitting.

CNN đại diện cho lớp biểu diễn học tương tác cục bộ không cần feature engineering. Tuy nhiên, một hạn chế căn bản là khái niệm "lân cận" trong véc tơ phẳng mang tính quy ước: hai cột allele liền kề có thể thuộc hai locus khác nhau, do đó phép tích chập không nhất thiết học được quan hệ sinh học thực sự giữa các peak trong cùng locus. Hạn chế này chính xác là những gì biểu diễn tập hợp phân cấp tìm cách khắc phục.

\begin{figure}[H]
\centering
\begin{tikzpicture}[
  box/.style={draw, rounded corners=4pt, align=center, font=\scriptsize, fill=blue!7},
  arr/.style={-Stealth, thick, gray}
]
\node[box, minimum width=1.8cm, minimum height=1.0cm, fill=orange!15] (inp)
  {\parbox{1.6cm}{\centering Véc tơ allele}};
\node[box, minimum width=2.6cm, minimum height=1.3cm, right=0.6cm of inp] (c1)
  {\parbox{2.4cm}{\centering Khối 1 \\ Conv 64, $k{=}3$ \\ BN + ReLU + MaxPool}};
\node[box, minimum width=2.6cm, minimum height=1.3cm, right=0.5cm of c1] (c2)
  {\parbox{2.4cm}{\centering Khối 2 \\ Conv 128, $k{=}3$ \\ BN + ReLU + MaxPool}};
\node[box, minimum width=2.6cm, minimum height=1.3cm, right=0.5cm of c2] (c3)
  {\parbox{2.4cm}{\centering Khối 3 \\ Conv 256, $k{=}3$ \\ BN + ReLU + AvgPool}};
\node[box, minimum width=2.2cm, minimum height=1.3cm, right=0.5cm of c3, fill=purple!12] (fc)
  {\parbox{2.0cm}{\centering FC \\ $2048{\to}256{\to}128$ \\ ReLU + Dropout}};
\node[box, minimum width=1.6cm, minimum height=1.0cm, right=0.5cm of fc, fill=red!10] (out)
  {\parbox{1.4cm}{\centering Softmax $\to$ NOC}};
\draw[arr] (inp) -- (c1);
\draw[arr] (c1)  -- (c2);
\draw[arr] (c2)  -- (c3);
\draw[arr] (c3)  -- (fc);
\draw[arr] (fc)  -- (out);
\end{tikzpicture}
\caption{Kiến trúc mô hình nền CNN một chiều. Ba khối tích chập học đặc trưng phân cấp từ chuỗi chiều cao allele; bộ phân loại fully connected hai lớp ánh xạ biểu diễn toàn cục sang phân phối xác suất năm lớp NOC.}
\label{fig:cnn_arch}
\end{figure}

\subsubsection{Mô hình nền Forensic Transformer}

Forensic Transformer tổ chức hồ sơ STR theo cấu trúc locus thay vì xếp phẳng toàn bộ allele (Hình~\ref{fig:ft_arch}). Đối với mỗi trong số 24 locus, véc tơ chiều cao allele của locus đó được chiếu riêng biệt qua một lớp tuyến tính tạo token 64 chiều. Thiết kế này cho phép mỗi locus có số allele khác nhau và không ràng buộc vào kích thước véc tơ toàn cục cố định. Chuỗi 24 token locus sau đó được xử lý qua ba lớp Transformer Encoder liên tiếp, mỗi lớp gồm multi-head self-attention và feed-forward network, cho phép mô hình học mối liên hệ xuyên locus trong toàn bộ profile. Các token đầu ra được gộp bằng mean pooling và đưa qua bộ phân loại gồm layer normalization và một lớp tuyến tính.

Forensic Transformer có khả năng nắm bắt tương tác toàn cục giữa các locus nhờ cơ chế self-attention — lợi thế rõ ràng so với CNN. Tuy nhiên, token của mỗi locus được tạo bằng phép chiếu tuyến tính đơn giản từ véc tơ allele phẳng của locus đó, nên thông tin về số peak, quan hệ chiều cao nội locus và hiệu ứng stutter vẫn bị nén mất trước khi vào Transformer. Đây là khác biệt căn bản so với DS-ST, nơi mỗi peak là một token độc lập được mã hóa với đầy đủ tám đặc trưng.

\begin{figure}[H]
\centering
\begin{tikzpicture}[
  box/.style={draw, rounded corners=4pt, align=center, font=\scriptsize, fill=blue!7},
  arr/.style={-Stealth, thick, gray}
]
\node[box, minimum width=1.8cm, minimum height=0.9cm, fill=orange!15] (l1)
  {\parbox{1.6cm}{\centering Locus 1 alleles}};
\node[box, minimum width=1.8cm, minimum height=0.9cm, below=0.28cm of l1, fill=orange!15] (l2)
  {\parbox{1.6cm}{\centering Locus 2 alleles}};
\node[below=0.12cm of l2, font=\normalsize, text=gray] (vd) {$\vdots$};
\node[box, minimum width=1.8cm, minimum height=0.9cm, below=0.12cm of vd, fill=orange!15] (l24)
  {\parbox{1.6cm}{\centering Locus 24 alleles}};

\node[box, minimum width=2.4cm, minimum height=2.8cm, right=0.9cm of l2] (lin)
  {\parbox{2.2cm}{\centering Linear \\ $d_{\ell} \to 64$ \\ (riêng từng locus)}};

\node[box, minimum width=2.6cm, minimum height=1.1cm, right=0.9cm of lin] (tr)
  {\parbox{2.4cm}{\centering $3\times$ Transformer Encoder \\ ($d{=}64$, $h{=}4$)}};

\node[box, minimum width=2.0cm, minimum height=1.0cm, right=0.7cm of tr] (mp)
  {\parbox{1.8cm}{\centering Mean Pooling 24 tokens}};

\node[box, minimum width=1.8cm, minimum height=1.0cm, right=0.7cm of mp, fill=red!10] (out)
  {\parbox{1.6cm}{\centering LN + Linear $\to$ NOC}};

\draw[arr] (l1.east)  -- ([xshift=0.3cm]l1.east) |- (lin.north west);
\draw[arr] (l2.east)  -- (lin.west);
\draw[arr] (l24.east) -- ([xshift=0.3cm]l24.east) |- (lin.south west);
\draw[arr] (lin) -- (tr);
\draw[arr] (tr)  -- (mp);
\draw[arr] (mp)  -- (out);
\end{tikzpicture}
\caption{Kiến trúc mô hình nền Forensic Transformer. Mỗi locus được chiếu độc lập thành token 64 chiều; ba lớp Transformer Encoder học tương tác xuyên locus qua self-attention; mean pooling trên 24 token tạo biểu diễn hồ sơ 64 chiều.}
\label{fig:ft_arch}
\end{figure}

\subsubsection{Mô hình nền TabPFN}

TabPFN là phương pháp meta-learning trong ngữ cảnh cho dữ liệu bảng, không yêu cầu bước tối ưu gradient tại thời điểm inference \cite{hollmann2023}. Trong quá trình tiền huấn luyện, mô hình được tối ưu trên hàng triệu bộ dữ liệu tổng hợp với phân phối tiên nghiệm Bayesian đa dạng, qua đó học các pattern phân loại phổ quát. Tại inference, toàn bộ tập huấn luyện được đưa vào context window của Transformer và dự đoán được sinh ra trong một forward pass duy nhất thông qua cơ chế in-context learning — về bản chất là Bayesian inference ngầm trên không gian hàm đã học.

Do giới hạn kích thước context window của mô hình, tập huấn luyện mỗi fold được lấy mẫu phân tầng tối đa 3.000 mẫu, nhỏ hơn đáng kể so với tập train đầy đủ của các mô hình khác. Mặc dù vậy, TabPFN xác nhận tiềm năng của in-context learning trên dữ liệu tabular STR và cung cấp điểm tham chiếu cho thấy các phương pháp tiền huấn luyện không cần cấu trúc domain-specific vẫn cạnh tranh tốt với tree boosting truyền thống.

\subsubsection{DS-ST — mô hình đề xuất}

DS-ST là kiến trúc trung tâm của nghiên cứu này, được thiết kế đặc thù để xử lý hồ sơ STR như một tập hợp phân cấp peak-locus. Toàn bộ chi tiết kiến trúc được trình bày trong Phần~\ref{sec:model}; phần này mô tả cấu hình sử dụng trong thực nghiệm.

Mô hình được khởi tạo từ checkpoint đã tiền huấn luyện trên nhiệm vụ contributor identification 45 lớp với 10.124 hồ sơ GlobalFiler. Encoder gồm hai khối ISAB++ với 32 inducing points trong không gian ẩn 128 chiều. Mỗi peak trong hồ sơ được biểu diễn bằng token tám chiều mã hóa đồng thời chỉ số locus, giá trị allele, logarithm chiều cao, cân bằng chiều cao trong locus, tỷ số stutter, xếp hạng giảm dần trong locus, số peak trong locus và chiều cao chuẩn hóa toàn hồ sơ. Periodic embedding ánh xạ giá trị allele liên tục sang không gian tuần hoàn tám chiều, cho phép mô hình nhận biết khoảng cách allele theo chu kỳ lặp đơn vị.

Fine-tune NOC được thực hiện qua hai giai đoạn. Giai đoạn đầu kéo dài năm epoch với backbone bị đóng băng, chỉ huấn luyện đầu NOC mới ở learning rate cao hơn; điều này ngăn biểu diễn tiền huấn luyện bị phá vỡ ngay từ đầu khi gradient đầu NOC còn lớn. Giai đoạn hai mở toàn bộ tham số với cosine learning rate scheduling trong tối đa 60 epoch và early stopping theo Macro F1 trên tập validation với patience bằng 12. Focal Loss với trọng số nghịch đảo tần suất lớp giảm ảnh hưởng của lớp NOC một chiếm hơn 80\% mẫu, buộc mô hình tập trung vào các lớp mixture phức tạp hơn.

\subsection{Kết quả benchmark}

\begin{table}[H]
\centering
\caption{Kết quả GroupKFold năm phần trên GlobalFiler. Năm mô hình đều dùng cùng phân chia group identity và cùng giao thức cân bằng lớp.}
\label{tab:benchmark}
\small
\begin{tabular}{lcc}
\toprule
\textbf{Mô hình} & \textbf{Macro F1} & \textbf{Micro F1} \\
\midrule
CNN một chiều       & \(0{,}8854\pm0{,}0100\) & \(0{,}9666\pm0{,}0051\) \\
Forensic Transformer & \(0{,}8846\pm0{,}0131\) & \(0{,}9624\pm0{,}0030\) \\
XGBoost             & \(0{,}9230\pm0{,}0209\) & \(0{,}9763\pm0{,}0082\) \\
TabPFN              & \(0{,}9297\pm0{,}0108\) & \(0{,}9787\pm0{,}0040\) \\
DS-ST  & \(\mathbf{0{,}9444\pm0{,}0042}\) & \(\mathbf{0{,}9822\pm0{,}0026}\) \\
\bottomrule
\end{tabular}
\end{table}

Bảng~\ref{tab:benchmark} trình bày kết quả GroupKFold năm phần của toàn bộ năm mô hình trên tập GlobalFiler. DS-ST dẫn đầu với Macro F1 $0{,}9444\pm0{,}0042$, vượt TabPFN $0{,}0147$ điểm và XGBoost $0{,}0214$ điểm tuyệt đối. Khoảng cách với CNN một chiều và Forensic Transformer lần lượt là $0{,}0590$ và $0{,}0598$ điểm Macro F1, phản ánh lợi ích của biểu diễn tập peak khi xử lý mixture phức tạp.

Một kết quả đáng chú ý là TabPFN ($0{,}9297\pm0{,}0108$) vượt XGBoost ($0{,}9230\pm0{,}0209$) dù không yêu cầu bước gradient descent tại inference time và chỉ sử dụng tối đa 3.000 mẫu huấn luyện mỗi fold do giới hạn context window. Điều này xác nhận rằng phương pháp Bayesian inference ngầm qua tiền huấn luyện in-context có thể cạnh tranh với tree boosting truyền thống trên dữ liệu tabular STR. Tuy nhiên, TabPFN vẫn thấp hơn DS-ST vì không khai thác cấu trúc tập peak theo locus.

DS-ST cũng là mô hình ổn định nhất giữa các fold với độ lệch chuẩn Macro F1 bằng $\pm0{,}0042$, thấp hơn TabPFN $\pm0{,}0108$, CNN $\pm0{,}0100$, Forensic Transformer $\pm0{,}0131$ và XGBoost $\pm0{,}0209$. Fold thấp nhất đạt $0{,}9370$, trong khi fold cao nhất đạt $0{,}9494$, tạo biên độ chỉ $0{,}0124$. Micro F1 của DS-ST bằng $0{,}9822\pm0{,}0026$; giá trị này được suy ra từ accuracy từng fold được log làm tròn ba chữ số, vì trong phân loại đơn nhãn Micro F1 bằng accuracy. Macro F1 vẫn là chỉ số chính do lớp NOC một chiếm hơn 80\% tổng mẫu.

\subsection{Phân tích theo lớp của DS-ST}

\begin{table}[H]
\centering
\caption{F1 theo lớp NOC của DS-ST trên từng fold GroupKFold 5. Tất cả giá trị được trích trực tiếp từ báo cáo phân loại.}
\label{tab:perclass}
\small
\resizebox{\textwidth}{!}{
\begin{tabular}{lccccccc}
\toprule
\textbf{Fold} & \textbf{NOC 1} & \textbf{NOC 2} & \textbf{NOC 3} & \textbf{NOC 4} & \textbf{NOC 5} & \textbf{Macro F1} & \textbf{Accuracy} \\
\midrule
1 & 0,996 & 0,939 & 0,967 & 0,934 & 0,911 & 0,9494 & 0,983 \\
2 & 0,995 & 0,967 & 0,974 & 0,917 & 0,863 & 0,9430 & 0,983 \\
3 & 0,992 & 0,894 & 0,920 & 0,934 & 0,945 & 0,9370 & 0,977 \\
4 & 0,995 & 0,917 & 0,949 & 0,958 & 0,913 & 0,9463 & 0,984 \\
5 & 0,994 & 0,943 & 0,948 & 0,895 & 0,953 & 0,9465 & 0,984 \\
\midrule
Trung bình & 0,994 & 0,932 & 0,952 & 0,928 & 0,917 & 0,9444 & 0,9822 \\
Độ lệch chuẩn & 0,001 & 0,025 & 0,019 & 0,021 & 0,032 & 0,0042 & 0,0026 \\
\bottomrule
\end{tabular}
}
\end{table}

Bảng~\ref{tab:perclass} trình bày F1 theo lớp của DS-ST trong năm fold. Lớp NOC một ổn định nhất với F1 từ $0{,}992$ đến $0{,}996$ và trung bình $0{,}994$. NOC ba đạt trung bình cao nhất trong các lớp mixture, bằng $0{,}952$, tiếp theo là NOC hai với $0{,}932$, NOC bốn với $0{,}928$ và NOC năm với $0{,}917$. NOC năm biến động mạnh nhất, từ $0{,}863$ ở fold hai đến $0{,}953$ ở fold năm; NOC hai cũng giảm xuống $0{,}894$ ở fold ba. Tuy vậy, mọi lớp mixture đều đạt F1 trung bình trên $0{,}91$, cho thấy mức Macro F1 cao không chỉ đến từ lớp NOC một chiếm đa số.

\begin{table}[H]
\centering
\caption{F1 theo lớp NOC của TabPFN trên từng fold GroupKFold 5. Tập huấn luyện được cap tối đa 3.000 mẫu/fold.}
\label{tab:perclass_tabpfn}
\small
\begin{tabular}{lccccc}
\toprule
\textbf{Fold} & \textbf{NOC 1} & \textbf{NOC 2} & \textbf{NOC 3} & \textbf{NOC 4} & \textbf{NOC 5} \\
\midrule
1 & 0,995 & 0,905 & 0,949 & 0,931 & 0,876 \\
2 & 0,992 & 0,895 & 0,953 & 0,875 & 0,833 \\
3 & 0,993 & 0,921 & 0,918 & 0,907 & 0,926 \\
4 & 0,995 & 0,900 & 0,966 & 0,952 & 0,898 \\
5 & 0,996 & 0,966 & 0,890 & 0,878 & 0,931 \\
\midrule
Trung bình & 0,994 & 0,917 & 0,935 & 0,909 & 0,893 \\
\midrule
Macro F1 & \multicolumn{5}{c}{$0{,}9297\pm0{,}0108$} \\
\bottomrule
\end{tabular}
\end{table}

Bảng~\ref{tab:perclass_tabpfn} cho thấy TabPFN duy trì hiệu năng tương đối ổn định trên NOC một và có khả năng cạnh tranh trên NOC ba. Tuy nhiên, DS-ST đạt F1 trung bình cao hơn TabPFN ở NOC hai ($0{,}932$ so với $0{,}917$), NOC ba ($0{,}952$ so với $0{,}935$), NOC bốn ($0{,}928$ so với $0{,}909$) và NOC năm ($0{,}917$ so với $0{,}893$). Như vậy, DS-ST dẫn trước trên cả bốn lớp mixture theo trung bình năm fold.

\subsection{Ablation study}

Ablation study sử dụng GroupKFold ba phần theo \texttt{group\_id}, cùng optimizer AdamW learning rate $\eta=3\times10^{-4}$, weight decay $10^{-4}$, batch size~64 và gradient clipping max\_norm=1{,}0. Mỗi fold thực hiện giai đoạn warmup đóng băng backbone năm epoch tại learning rate $3\eta=9\times10^{-4}$, tiếp theo huấn luyện toàn mạng tối đa 50 epoch với early stopping patience~10 và CosineAnnealingLR:
\begin{equation}
\eta_t = \eta_{\min} + \tfrac{1}{2}\!\left(\eta - \eta_{\min}\right)\!\left(1 + \cos\tfrac{\pi t}{T_{\max}}\right),\quad T_{\max}=50,\;\eta_{\min}=10^{-6}.
\end{equation}
Tập huấn luyện mỗi fold được oversample để cân bằng số mẫu giữa năm lớp NOC trước khi đưa vào DataLoader. Hai biến thể sau được thiết kế để tách riêng đóng góp của từng thành phần so với cấu hình DS-ST đầy đủ.

\paragraph{Biến thể không tiền huấn luyện (\texttt{no\_pretrain}).} Backbone được khởi tạo ngẫu nhiên thay vì nạp checkpoint đã tiền huấn luyện trên DS-ST. Toàn bộ kiến trúc ISAB++, non-competitive sigmoid attention, PLR embedding, token tám chiều, set-of-set, feasibility filter và Focal Loss giữ nguyên. Biến thể này định lượng đóng góp độc lập của nhiệm vụ tiền huấn luyện contributor identification 45~lớp.

\paragraph{Biến thể không enrichment (\texttt{no\_enrich}).} Backbone nạp checkpoint DS-ST tiền huấn luyện nhưng token bị rút gọn về ba trường cơ sở với năm chiều còn lại đặt bằng không:
\begin{equation}
\bm{x}_p^{\mathrm{(plain)}} = \bigl[\,l,\; a,\; \log(1+h),\; 0,\; 0,\; 0,\; 0,\; 0\,\bigr].
\end{equation}
Cấu hình đầy đủ dùng token tám chiều:
\begin{equation}
\bm{x}_p = \Bigl[\,l,\; a,\; \log(1+h),\; \mathrm{Hb},\; \mathrm{SR},\; \tfrac{\mathrm{rank}}{n_l-1},\; \tfrac{n_l}{12},\; \tfrac{h}{h_{\max}}\,\Bigr],
\end{equation}
trong đó $\mathrm{Hb}=h/\!\sum_{p'\in l}h_{p'}$ là cân bằng chiều cao trong locus, $\mathrm{SR}=h/(h_{a+1}+\varepsilon)$ là tỷ số với peak cha stutter (bằng không khi không tồn tại peak cha), và rank tính giảm dần trong locus (0 = cao nhất). Biến thể này kiểm tra đóng góp của năm trường enrichment bổ sung so với đặc trưng ba chiều cơ sở.


\begin{table}[H]
\centering
\caption{Ablation GroupKFold ba phần. Cấu hình DS-ST đầy đủ là baseline trong cùng protocol. Dấu~†~chỉ kết quả cần chạy bổ sung.}
\label{tab:ablation}
\small
\begin{tabular}{lccc}
\toprule
\textbf{Cấu hình} & \textbf{Macro F1} & \textbf{Micro F1} & \textbf{Thay đổi Macro} \\
\midrule
DS-ST đầy đủ & \(0,9122\pm0,0089\) & \(0,9709\pm0,0020\) & \(0,0000\) \\
Không tiền huấn luyện (\texttt{no\_pretrain}) & \(0,8918\pm0,0186\) & \(0,9657\pm0,0049\) & \(-0,0204\) \\
Không enrichment, token 3 chiều (\texttt{no\_enrich})$^\dagger$ & {--} & {--} & {--} \\
Checkpoint và pipeline cũ & \(0,8430\pm0,0156\) & \(0,9361\pm0,0074\) & \(-0,0692\) \\
\bottomrule
\end{tabular}
\end{table}

Kết quả không tiền huấn luyện cho thấy Macro F1 giảm 2,04 điểm phần trăm và độ lệch chuẩn tăng từ 0,0089 lên 0,0186. Điều này cho thấy tiền huấn luyện không chỉ cải thiện trung bình mà còn tăng tính ổn định giữa các fold. Thí nghiệm pipeline cũ đồng thời thay thế nhiều thành phần: checkpoint tiền huấn luyện thế hệ trước (huấn luyện với encoder ISAB cơ sở không có cải tiến inducing point), encoder ISAB cơ sở, embedding raw không có PLR, token ba chiều, không phân tách private/shared peak và không feasibility filter. Macro F1 giảm $6{,}92$ điểm phần trăm, nhưng vì nhiều thành phần thay đổi đồng thời, kết quả chỉ chứng minh toàn bộ pipeline hiện đại vượt trội so với thế hệ trước mà chưa tách được đóng góp riêng của PLR embedding, token enrichment hoặc ISAB++. Ablation đơn biến \texttt{no\_enrich} được thiết kế để giải quyết câu hỏi này và kết quả sẽ được cập nhật sau khi hoàn tất.

\subsection{Kết quả nhận dạng contributor — Bài toán 1 (DS-ST)}

Bài toán 1 yêu cầu mô hình xác định tập contributor từ panel tham chiếu $C=45$ donor, đồng thời từ chối hồ sơ open-set (chứa contributor ngoài panel). Tập kiểm thử gồm 1.358 hồ sơ (1.125 NOC\,=\,1, 44 NOC\,=\,2, 96 NOC\,=\,3, 45 NOC\,=\,4, 48 NOC\,=\,5). Baseline LR, XGBoost và CNN được huấn luyện trên dữ liệu thật; DS-ST và Hybrid (ensemble DS-ST với LR) được huấn luyện trên 50k hồ sơ GF29cycles tổng hợp.

\begin{table}[H]
\centering
\caption{Kết quả nhận dạng contributor trên tập kiểm thử. Exact Match (EM): toàn bộ tập contributor dự đoán đúng. Oracle EM: dùng NOC thật để chọn top $k$ donor (đánh giá chất lượng xếp hạng độc lập với lỗi đếm). Reject AUROC: phân biệt closed-set và open-set. Ký hiệu \textsuperscript{†}: mô hình dùng 50k dữ liệu GF29cycles tổng hợp.}
\label{tab:donor_id}
\small
\begin{tabular}{lccccc}
\toprule
\textbf{Mô hình} & \textbf{Macro F1} & \textbf{Exact Match} & \textbf{Oracle EM} & \textbf{Jaccard} & \textbf{Reject AUROC} \\
\midrule
LR (baseline)          & 0,935 & 0,847 & ---   & 0,911 & --- \\
XGBoost (baseline)     & 0,941 & 0,845 & ---   & 0,921 & --- \\
CNN (baseline)         & 0,899 & 0,787 & ---   & 0,872 & --- \\
DS-ST\textsuperscript{†} & 0,929 & 0,854 & 0,903 & 0,926 & 0,9991 \\
Hybrid (DS-ST + LR)\textsuperscript{†} & \textbf{0,980} & \textbf{0,954} & \textbf{0,975} & \textbf{0,980} & 0,9956 \\
\bottomrule
\end{tabular}
\end{table}

Bảng~\ref{tab:donor_id} cho thấy mô hình Hybrid đạt Exact Match 0,954, vượt DS-ST đơn lẻ (0,854) 10 điểm phần trăm và vượt tất cả baseline. Oracle EM của Hybrid đạt 0,975, nghĩa là nếu biết đúng NOC thì mô hình xếp hạng đúng 97,5\% hồ sơ. Khoảng cách giữa Oracle EM (0,975) và EM thực tế (0,954) phản ánh đóng góp của bước NOC estimation trong pipeline two-stage.

Cả DS-ST và Hybrid đều có Reject AUROC trên 0,995, cho thấy khả năng phân biệt rất tốt giữa hồ sơ closed-set (mọi contributor trong panel) và open-set (có contributor ngoài panel). Baseline LR, XGBoost và CNN không được trang bị cơ chế từ chối open-set nên không báo cáo chỉ số này.

\begin{table}[H]
\centering
\caption{Exact Match theo NOC trên tập kiểm thử (Bài toán 1, $n_{\mathrm{total}}=1358$). Hybrid dùng giải mã two-stage+pgNOC; DS-ST dùng giải mã post-hoc. Cỡ mẫu: NOC\,1: 1125, NOC\,2: 44, NOC\,3: 96, NOC\,4: 45, NOC\,5: 48.}
\label{tab:donor_id_pernoc}
\footnotesize
\begin{adjustbox}{max width=\textwidth}
\begin{tabular}{lcccccc}
\toprule
\textbf{Mô hình} & \textbf{NOC\,=\,1} & \textbf{NOC\,=\,2} & \textbf{NOC\,=\,3} & \textbf{NOC\,=\,4} & \textbf{NOC\,=\,5} \\
\midrule
LR (baseline)   & 0,928 & 0,636 & 0,750 & 0,111 & 0,021 \\
XGBoost         & 0,943 & 0,364 & 0,698 & 0,000 & 0,083 \\
CNN             & 0,941 & 0,227 & 0,000 & 0,000 & 0,000 \\
DS-ST\textsuperscript{†}  & \textbf{0,973} & 0,045 & 0,646 & 0,000 & 0,000 \\
Hybrid\textsuperscript{†} & 0,996 & 0,795 & \textbf{0,917} & \textbf{0,733} & \textbf{0,396} \\
\bottomrule
\end{tabular}
\end{adjustbox}
\end{table}

Bảng~\ref{tab:donor_id_pernoc} làm rõ điểm yếu cấu trúc của từng mô hình. LR đạt EM cao nhất ở NOC\,=\,2 (0,636) trong nhóm baseline, nhưng sụp đổ ở NOC\,=\,4 (0,111) và NOC\,=\,5 (0,021). XGBoost tốt hơn ở NOC\,=\,1 và NOC\,=\,3 nhưng không phát hiện được mixture bốn và năm người (EM = 0). CNN hoàn toàn thất bại từ NOC\,=\,3 trở lên (EM = 0). DS-ST nhận dạng tốt NOC\,=\,1 (0,973) nhưng gặp khó ở NOC\,=\,2 (0,045) — bằng chứng cho thấy biểu diễn slot bị nhầm lẫn giữa hai donor chia sẻ allele khi không có tín hiệu LR bổ trợ. Hybrid giải quyết vấn đề này bằng cách kết hợp xác suất slot của DS-ST với xếp hạng donor từ LR, cho phép NOC\,=\,4 đạt EM 0,733 và NOC\,=\,5 đạt 0,396 — mức cao nhất trong tất cả mô hình.

\subsection{So sánh với nghiên cứu trước}

\begin{table}[H]
\centering
\caption{Đối chiếu các nghiên cứu NOC. Chỉ số không được xem là so sánh trực tiếp vì dữ liệu, số lớp và protocol khác nhau.}
\label{tab:related_compare}
\small
\begin{tabularx}{\textwidth}{P{2.5cm}P{1.7cm}P{2.5cm}P{2.3cm}Y}
\toprule
\textbf{Nghiên cứu} & \textbf{Miền NOC} & \textbf{Dữ liệu} & \textbf{Kết quả báo cáo} & \textbf{Khác biệt chính} \\
\midrule
PACE \cite{marciano2017} & 1 đến 4 & 1.405 profile, 20 donor, 15 locus & khoảng 97\% test accuracy & feature engineering và Identifiler \\
RFC19 \cite{benschop2019} & 1 đến 5 & 590 PPF6C profile, 1.174 donor & khoảng 83,3\% accuracy & random forest 19 feature \\
Decision tree \cite{kruijver2021} & 1 đến 5 & 766 PROVEDIt profile & khoảng 0,77 đến 0,85 accuracy & minh bạch, phụ thuộc lọc artefact \\
TAWSEEM \cite{alamoudi2022} & 1 đến 5 & 6.000 profile, bốn multiplex & 97\% accuracy & MLP trên khoảng 6.000 feature \\
deepNoC \cite{taylor2024deepnoc} & 1 đến 10 & 100.000 profile mô phỏng & 89\% accuracy & simulator và fine tune phòng lab \\
TabPFN (thử nghiệm nội bộ) & 1 đến 5 & 10.124 GF profile, GroupKFold 5, cap 3.000/fold & Macro F1 $0{,}9297\pm0{,}0108$; Micro F1 $0{,}9787\pm0{,}0040$ & in-context learning; không cần fine-tune; đặc trưng phẳng \\
Nghiên cứu này — DS-ST & 1 đến 5 & 10.124 GlobalFiler profile, 3.418 group & Macro F1 $0{,}9444\pm0{,}0042$; Micro F1 $0{,}9822\pm0{,}0026$ & GroupKFold, token 8 chiều, tiền huấn luyện contributor \\
\bottomrule
\end{tabularx}
\end{table}

Bảng~\ref{tab:related_compare} cho thấy accuracy cao không tự động đồng nghĩa phương pháp tốt hơn trong mọi môi trường. TAWSEEM dùng nhiều multiplex và véc tơ lớn, deepNoC dùng dữ liệu mô phỏng cho mười lớp, PACE chỉ đến bốn contributor, trong khi nghiên cứu hiện tại nhấn mạnh GroupKFold theo mixture identity và Macro F1 trên dữ liệu mất cân bằng. Một so sánh công bằng trong tương lai cần cùng dataset version, cùng threshold artefact, cùng split và cùng metric.

\subsection{Khả năng tổng quát hóa khác kit}

Identifiler Plus và PowerPlex 16 có lần lượt 3.160 và 1.022 hồ sơ. Hai tập nhỏ hơn GlobalFiler và có phân phối mất cân bằng tương tự. Kiến trúc phân cấp xử lý khác biệt kit bằng mask locus và giữ locus mapping chuẩn, thay vì sắp lại locus theo từng tập. Trong thử nghiệm cross kit từ đầu, backbone dùng cùng hidden size và loss nhưng không nạp checkpoint GlobalFiler khi locus không tương thích. Các run cross kit hoàn chỉnh cần được lưu dưới dạng artifact trước khi báo cáo con số chính thức; do đó bản paper này chỉ mô tả protocol và thống kê dữ liệu, không suy diễn metric chưa có bằng chứng.

\subsection{Thảo luận}

Kết quả benchmark cho thấy hai bức tranh rõ nét. Thứ nhất, các phương pháp tabular truyền thống như XGBoost và mô hình tiền huấn luyện in-context như TabPFN vẫn rất cạnh tranh trên dữ liệu phẳng; TabPFN vượt XGBoost dù chỉ dùng tối đa 3.000 mẫu huấn luyện mỗi fold, phù hợp với kết quả của Hollmann và cộng sự trên dữ liệu tabular quy mô nhỏ \cite{hollmann2023}. Thứ hai, DS-ST đạt Macro F1 $0{,}9444\pm0{,}0042$, cao hơn cả hai baseline và đồng thời có độ lệch chuẩn thấp nhất. Lợi thế có thể đến từ tổ hợp tiền huấn luyện contributor, token tám chiều, periodic embedding và ISAB++; ablation xác nhận giá trị của toàn pipeline nhưng chưa phân bổ riêng mức tăng cho từng thành phần.

Kiến trúc phân cấp peak-locus-profile có ưu điểm diễn giải cấu trúc rõ hơn DS-ST nhánh NoC. Mỗi cấp xử lý tương ứng với một đơn vị pháp y: peak, locus và profile. Stutter bias attention cho phép kiểm tra attention theo khoảng cách allele; auxiliary allele count cho phép kiểm tra locus token; ba đầu NoC cho phép quan sát khi phân loại, ordinal và regression bất đồng. Bất đồng giữa ba đầu có thể trở thành tín hiệu uncertainty trong tương lai.

Nhánh contributor identification biến donor thành slot có điều kiện genotype. Đây là phép chuyển khái niệm từ object centric learning sang forensic mixture. Sự tương đồng về cấu trúc là mỗi input gồm nhiều quan sát chồng lấn và cần tách về nhiều nguồn ẩn. Tuy nhiên, donor không phải vật thể ảnh và allele sharing tuân theo quy luật di truyền, nên các module CoSA, GSA, MESH và AdaSlot chỉ được xem là nguồn cảm hứng kiến trúc; validation pháp y vẫn phải dựa trên dữ liệu mixture thực.

Một hạn chế quan trọng là các protocol hiện tại chưa hoàn toàn thống nhất. Benchmark chính sử dụng năm fold, ablation dùng ba fold, còn contributor reranking mới có kết quả seed 42. Paper phân tách rõ các protocol để tránh trộn số liệu. Phiên bản tiếp theo cần chạy toàn bộ ablation năm fold, nhiều seed cho contributor ID và một benchmark cross kit frozen hoặc fine tune có artifact đầy đủ.

\subsection{Các đóng góp chính của bài báo}

Kết quả thực nghiệm hỗ trợ bốn đóng góp. Biểu diễn tập hợp giúp DS-ST đạt Macro F1 $0{,}9444\pm0{,}0042$, cao hơn XGBoost $0{,}0214$ điểm và TabPFN $0{,}0147$ điểm dưới cùng GroupKFold năm phần trên 10.124 hồ sơ GlobalFiler. Năm fold lần lượt đạt $0{,}9494$, $0{,}9430$, $0{,}9370$, $0{,}9463$ và $0{,}9465$, cho thấy độ ổn định cao. Tiền huấn luyện contributor giúp Macro F1 tăng $0{,}0204$ trong ablation ba phần. Toàn bộ pipeline giàu đặc trưng và encoder cải tiến tăng $0{,}0692$ so với checkpoint và pipeline cũ.

\section{Kết luận và hướng phát triển}

\subsection{Kết luận}

Nghiên cứu đã xây dựng một khung học sâu cho hai nhiệm vụ trung tâm của phân tích mixture DNA STR: dự đoán Number of Contributors và xác định contributor trong panel tham chiếu. Điểm cốt lõi là xem hồ sơ như một tập hợp phân cấp. Peak Embedder học đặc trưng cục bộ, Locus Encoder học quan hệ peak và stutter, Deep Sets tạo locus token bất biến hoán vị, Cross Locus Transformer học ngữ cảnh toàn hồ sơ, còn Profile Pool cung cấp véc tơ dùng chung cho nhiều đầu nhiệm vụ.

Trên 10.124 hồ sơ GlobalFiler thuộc 3.418 mixture group với GroupKFold năm phần, DS-ST đạt Macro F1 $0{,}9444\pm0{,}0042$, vượt TabPFN ($0{,}9297\pm0{,}0108$), XGBoost ($0{,}9230\pm0{,}0209$), CNN một chiều ($0{,}8854\pm0{,}0100$) và Forensic Transformer ($0{,}8846\pm0{,}0131$). Ablation cho thấy bỏ tiền huấn luyện làm giảm Macro F1 $0{,}0204$ và tăng độ lệch chuẩn từ $0{,}0089$ lên $0{,}0186$; quay về checkpoint và pipeline cũ làm giảm $0{,}0692$. Kết quả TabPFN ($0{,}9297$) vượt XGBoost trên cùng đặc trưng phẳng, xác nhận tiềm năng của phương pháp in-context learning, nhưng DS-ST vẫn đạt hiệu năng tốt nhất và ổn định nhất trên mixture phức tạp.

Đối với contributor identification, bộ giải mã genotype conditioned slot cung cấp một cách diễn giải trực tiếp: mỗi slot ứng với một donor và ma trận optimal transport thể hiện peak nào hỗ trợ donor nào. Phi reranking seed 42 làm oracle exact match NOC năm tăng từ 0,831 lên 0,901, nhưng cần đánh giá nhiều seed trước khi coi mức tăng là ổn định.

\subsection{Hướng phát triển}

Hướng phát triển gần nhất là hoàn thành ablation đơn biến cho từng thành phần, gồm bỏ enrichment (\texttt{no\_enrich}), bỏ stutter bias, bỏ CORN, bỏ auxiliary allele count và thay profile pooling. Tất cả cấu hình cần dùng cùng GroupKFold năm phần và cùng seed list để có thể tính confidence interval.

Hướng thứ hai là huấn luyện đa nhiệm end to end. NOC head và contributor head hiện chia sẻ backbone về mặt kiến trúc nhưng có thể được tối ưu theo giai đoạn. Một loss consistency có thể ép tổng gate contributor gần kỳ vọng NOC:
\begin{equation}
\mathcal{L}_{\mathrm{cons}}=
\left|
\sum_{c=1}^{C}\sigma(e_c)
-\sum_{k=1}^{5}k\,p_k^{\mathrm{noc}}
\right|.
\end{equation}
Thiết kế này cho phép hai nhánh hỗ trợ nhau nhưng cần kiểm soát để lỗi đếm không làm hỏng ranking donor.

Hướng thứ ba là đánh giá cross kit có kiểm soát trên Identifiler Plus và PowerPlex 16. Cần so sánh train from scratch, frozen GlobalFiler encoder, fine tune toàn bộ và adapter theo kit. Locus không tồn tại phải giữ mask, còn locus chung phải giữ cùng chỉ số. Báo cáo cần có Macro F1, calibration, per class F1 và performance theo mixture ratio.

Hướng thứ tư là tích hợp allele calling từ electropherogram thô. DNANet cho thấy U Net có thể đạt F1 0,962 trên allele donor thật \cite{dewit2025}. Một pipeline kết hợp DNANet với kiến trúc phân cấp peak-locus-profile có thể truyền uncertainty ở mức scan point tới peak và profile, thay vì coi peak extraction là hoàn toàn chính xác.

Hướng cuối cùng là phát triển explainability phù hợp pháp y. Attention map không tự động là explanation. Cần kết hợp counterfactual profile hợp lệ như Veldhuis và cộng sự, kiểm tra sensitivity theo peak dropout, stutter và thay đổi mixture proportion, đồng thời đánh giá explanation với chuyên gia \cite{veldhuis2022}. Mục tiêu cuối không chỉ là tăng metric mà là tạo hệ thống có thể kiểm chứng, tái lập và sử dụng như công cụ hỗ trợ thay vì thay thế phán đoán chuyên môn.

\newpage
\begin{thebibliography}{99}

\bibitem{jeffreys1985}
Jeffreys, A. J., Wilson, V., \& Thein, S. L. (1985).
Individual-specific fingerprints of human DNA.
\textit{Nature}, 316, 76--79.
\url{https://doi.org/10.1038/316076a0}

\bibitem{nist2024}
National Institute of Standards and Technology. (2024).
\textit{DNA Mixture Interpretation: A Scientific Foundation Review}.
NIST IR 8351.
\url{https://doi.org/10.6028/NIST.IR.8351}

\bibitem{provedit2018}
Alfonse, L. E., Garrett, A. D., Lun, D. S., Duffy, K. R., \& Grgicak, C. M. (2018).
A large-scale dataset of single and mixed-source short tandem repeat profiles to inform human identification strategies: PROVEDIt.
\textit{Forensic Science International: Genetics}, 32, 62--70.
\url{https://doi.org/10.1016/j.fsigen.2017.10.006}

\bibitem{proveditweb}
Laboratory for Forensic Technology Development and Integration. (2026).
PROVEDIt Database.
\url{https://lftdi.camden.rutgers.edu/provedit/}

\bibitem{haned2011}
Haned, H., Pène, L., Lobry, J. R., Dufour, A. B., \& Pontier, D. (2011).
Estimating the number of contributors to forensic DNA mixtures: Does maximum likelihood perform better than maximum allele count?
\textit{Journal of Forensic Sciences}, 56(1), 23--28.
\url{https://doi.org/10.1111/j.1556-4029.2010.01550.x}

\bibitem{alfonse2017}
Alfonse, L. E., Tejada, G., Swaminathan, H., Lun, D. S., \& Grgicak, C. M. (2017).
Inferring the number of contributors to complex DNA mixtures using three methods: Exploring the limits of low-template DNA interpretation.
\textit{Journal of Forensic Sciences}, 62(2), 308--316.
\url{https://doi.org/10.1111/1556-4029.13284}

\bibitem{grgicak2020}
Grgicak, C. M., Karkar, S., Yearwood-Garcia, X., Alfonse, L. E., Duffy, K. R., \& Lun, D. S. (2020).
A large-scale validation of NOCIt's a posteriori probability of the number of contributors and its integration into forensic interpretation pipelines.
\textit{Forensic Science International: Genetics}, 47, 102296.
\url{https://doi.org/10.1016/j.fsigen.2020.102296}

\bibitem{marciano2017}
Marciano, M. A., \& Adelman, J. D. (2017).
PACE: Probabilistic Assessment for Contributor Estimation, a machine learning-based assessment of the number of contributors in DNA mixtures.
\textit{Forensic Science International: Genetics}, 27, 82--91.
\url{https://doi.org/10.1016/j.fsigen.2016.11.006}

\bibitem{benschop2019}
Benschop, C. C. G., van der Linden, J., Hoogenboom, J., Ypma, R. J. F., \& Haned, H. (2019).
Automated estimation of the number of contributors in autosomal short tandem repeat profiles using a machine learning approach.
\textit{Forensic Science International: Genetics}, 43, 102150.
\url{https://doi.org/10.1016/j.fsigen.2019.102150}

\bibitem{kruijver2021}
Kruijver, M., Kelly, H., Cheng, K., Lin, M. H., Morawitz, J., Russell, L., Buckleton, J., \& Bright, J. A. (2021).
Estimating the number of contributors to a DNA profile using decision trees.
\textit{Forensic Science International: Genetics}, 50, 102407.
\url{https://doi.org/10.1016/j.fsigen.2020.102407}

\bibitem{alamoudi2022}
Alamoudi, E., Mehmood, R., Albeshri, A., \& Gojobori, T. (2022).
TAWSEEM: A deep-learning-based tool for estimating the number of unknown contributors in DNA profiling.
\textit{Electronics}, 11(4), 548.
\url{https://doi.org/10.3390/electronics11040548}

\bibitem{veldhuis2022}
Veldhuis, M. S., Ariëns, S., Ypma, R. J. F., Abeel, T. E. P. M. F., \& Benschop, C. C. G. (2022).
Explainable artificial intelligence in forensics: Realistic explanations for number of contributor predictions of DNA profiles.
\textit{Forensic Science International: Genetics}, 56, 102632.
\url{https://doi.org/10.1016/j.fsigen.2021.102632}

\bibitem{taylor2024deepnoc}
Taylor, D., \& Humphries, M. A. (2024).
deepNoC: A deep learning system to assign the number of contributors to a short tandem repeat DNA profile.
\textit{arXiv preprint arXiv:2412.09803}.
\url{https://doi.org/10.48550/arXiv.2412.09803}

\bibitem{barash2024}
Barash, M., McNevin, D., Fedorenko, V., \& Giverts, P. (2024).
Machine learning applications in forensic DNA profiling: A critical review.
\textit{Forensic Science International: Genetics}, 69, 102994.
\url{https://doi.org/10.1016/j.fsigen.2023.102994}

\bibitem{phan2021}
Phan, N. N., Chattopadhyay, A., Lee, T. T., Yin, H. I., Lu, T. P., Lai, L. C., Hwa, H. L., Tsai, M. H., \& Chuang, E. Y. (2021).
High-performance deep learning pipeline predicts individuals in mixtures of DNA using sequencing data.
\textit{Briefings in Bioinformatics}, 22(6), bbab283.
\url{https://doi.org/10.1093/bib/bbab283}

\bibitem{dewit2025}
de Wit, A. K. J. G., Wagenaar, C. D., Janssen, N. A. C., Hoegen, B., van de Wetering, J., Hoofs, H., Ariëns, S., Benschop, C. C. G., \& Ypma, R. J. F. (2026).
Making AI accessible for forensic DNA profile analysis.
\textit{Forensic Science International: Genetics}, 81, 103345. Published online August 19, 2025.
\url{https://doi.org/10.1016/j.fsigen.2025.103345}

\bibitem{bleka2016}
Bleka, Ø., Storvik, G., \& Gill, P. (2016).
EuroForMix: An open source software based on a continuous model to evaluate STR DNA profiles from a mixture of contributors with artefacts.
\textit{Forensic Science International: Genetics}, 21, 35--44.
\url{https://doi.org/10.1016/j.fsigen.2015.11.008}

\bibitem{bright2016}
Bright, J. A., Taylor, D., McGovern, C., Cooper, S., Russell, L., Abarno, D., \& Buckleton, J. (2016).
Developmental validation of STRmix, expert software for the interpretation of forensic DNA profiles.
\textit{Forensic Science International: Genetics}, 23, 226--239.
\url{https://doi.org/10.1016/j.fsigen.2016.05.007}

\bibitem{coble2019}
Coble, M. D., \& Bright, J. A. (2019).
Probabilistic genotyping software: An overview.
\textit{Forensic Science International: Genetics}, 38, 219--224.
\url{https://doi.org/10.1016/j.fsigen.2018.11.009}

\bibitem{vaswani2017}
Vaswani, A., Shazeer, N., Parmar, N., Uszkoreit, J., Jones, L., Gomez, A. N., Kaiser, Ł., \& Polosukhin, I. (2017).
Attention is all you need.
\textit{Advances in Neural Information Processing Systems}, 30, 5998--6008.
\url{https://proceedings.neurips.cc/paper/7181-attention-is-all-you-need}

\bibitem{zaheer2017}
Zaheer, M., Kottur, S., Ravanbakhsh, S., Póczos, B., Salakhutdinov, R., \& Smola, A. J. (2017).
Deep Sets.
\textit{Advances in Neural Information Processing Systems}, 30, 3391--3401.
\url{https://papers.nips.cc/paper/6931-deep-sets}

\bibitem{lee2019}
Lee, J., Lee, Y., Kim, J., Kosiorek, A., Choi, S., \& Teh, Y. W. (2019).
Set Transformer: A framework for attention-based permutation-invariant neural networks.
\textit{Proceedings of the 36th International Conference on Machine Learning}, 97, 3744--3753.
\url{https://proceedings.mlr.press/v97/lee19d.html}

\bibitem{zhang2023setnorm}
Zhang, L., Tozzo, V., Higgins, J. M., \& Ranganath, R. (2022).
Set norm and equivariant skip connections: Putting the deep in deep sets.
\textit{Proceedings of the 39th International Conference on Machine Learning}, 162, 26559--26574.
\url{https://proceedings.mlr.press/v162/zhang22ac.html}

\bibitem{gorishniy2022}
Gorishniy, Y., Rubachev, I., \& Babenko, A. (2022).
On embeddings for numerical features in tabular deep learning.
\textit{Advances in Neural Information Processing Systems}, 35, 24991--25004.
\url{https://proceedings.neurips.cc/paper_files/paper/2022/hash/9e9f0ffc3d836836ca96cbf8fe14b105-Abstract-Conference.html}

\bibitem{hendrycks2016}
Hendrycks, D., \& Gimpel, K. (2016).
Gaussian error linear units.
\textit{arXiv preprint arXiv:1606.08415}.
\url{https://doi.org/10.48550/arXiv.1606.08415}

\bibitem{lin2017}
Lin, T. Y., Goyal, P., Girshick, R., He, K., \& Dollár, P. (2017).
Focal Loss for Dense Object Detection.
\textit{Proceedings of the IEEE International Conference on Computer Vision}, 2980--2988.
\url{https://doi.org/10.1109/ICCV.2017.324}

\bibitem{shi2023corn}
Shi, X., Cao, W., \& Raschka, S. (2023).
Deep neural networks for rank-consistent ordinal regression based on conditional probabilities.
\textit{Pattern Analysis and Applications}, 26, 941--955.
\url{https://doi.org/10.1007/s10044-023-01181-9}

\bibitem{chen2016}
Chen, T., \& Guestrin, C. (2016).
XGBoost: A scalable tree boosting system.
\textit{Proceedings of the 22nd ACM SIGKDD International Conference on Knowledge Discovery and Data Mining}, 785--794.
\url{https://doi.org/10.1145/2939672.2939785}

\bibitem{hollmann2023}
Hollmann, N., Müller, S., Eggensperger, K., \& Hutter, F. (2023).
TabPFN: A Transformer that solves small tabular classification problems in a second.
\textit{International Conference on Learning Representations}.
\url{https://openreview.net/forum?id=eu9fVjVasr4}

\bibitem{locatello2020slot}
Locatello, F., Weissenborn, D., Unterthiner, T., Mahendran, A., Heigold, G., Uszkoreit, J., Dosovitskiy, A., \& Kipf, T. (2020).
Object-centric learning with Slot Attention.
\textit{Advances in Neural Information Processing Systems}, 33, 11525--11538.
\url{https://proceedings.neurips.cc/paper/2020/hash/8511df98c02ab60aea1b2356c013bc0f-Abstract.html}

\bibitem{ma2024cosa}
Kori, A., Locatello, F., De Sousa Ribeiro, F., Toni, F., \& Glocker, B. (2024).
Grounded object-centric learning.
\textit{International Conference on Learning Representations}.
\url{https://openreview.net/forum?id=pBxeZ6pVUD}

\bibitem{lee2024gsa}
Lee, M., Cho, S., Lee, D., Park, C., Lee, J., \& Lee, S. (2024).
Guided Slot Attention for unsupervised video object segmentation.
\textit{Proceedings of the IEEE/CVF Conference on Computer Vision and Pattern Recognition}.
\url{https://doi.org/10.48550/arXiv.2303.08314}

\bibitem{zhang2023mesh}
Zhang, Y., Zhang, D. W., Lacoste-Julien, S., Burghouts, G. J., \& Snoek, C. G. M. (2023).
Unlocking Slot Attention by changing optimal transport costs.
\textit{Proceedings of the 40th International Conference on Machine Learning}, 202, 41931--41951.
\url{https://proceedings.mlr.press/v202/zhang23ba.html}

\bibitem{fan2024adaslot}
Fan, K., Bai, Z., Xiao, T., He, T., Horn, M., Fu, Y., Locatello, F., \& Zhang, Z. (2024).
Adaptive Slot Attention: Object discovery with dynamic slot number.
\textit{Proceedings of the IEEE/CVF Conference on Computer Vision and Pattern Recognition}, 23062--23071.
\url{https://openaccess.thecvf.com/content/CVPR2024/html/Fan_Adaptive_Slot_Attention_Object_Discovery_with_Dynamic_Slot_Number_CVPR_2024_paper.html}

\end{thebibliography}

\end{document}
